\documentclass[a4paper,12pt,twoside]{memoir}

% Castellano
\usepackage[spanish,es-tabla]{babel}
\selectlanguage{spanish}
\usepackage[utf8]{inputenc}
\usepackage[T1]{fontenc}
\usepackage{lmodern} % scalable font
\usepackage{microtype}
\usepackage{placeins}

\RequirePackage{booktabs}
\RequirePackage[table]{xcolor}
\RequirePackage{xtab}
\RequirePackage{multirow}

% Links
\PassOptionsToPackage{hyphens}{url}\usepackage[colorlinks]{hyperref}
\hypersetup{
	allcolors = {red}
}

% Ecuaciones
\usepackage{amsmath}

% Rutas de fichero / paquete
\newcommand{\ruta}[1]{{\sffamily #1}}

% Párrafos
\nonzeroparskip

% Huérfanas y viudas
\widowpenalty100000
\clubpenalty100000

% Evitar solapes en el header
\nouppercaseheads

% Imagenes
\usepackage{graphicx}
\newcommand{\imagen}[2]{
	\begin{figure}[!h]
		\centering
		\includegraphics[width=0.9\textwidth]{#1}
		\caption{#2}\label{fig:#1}
	\end{figure}
	\FloatBarrier
}

\newcommand{\imagenflotante}[2]{
	\begin{figure}%[!h]
		\centering
		\includegraphics[width=0.9\textwidth]{#1}
		\caption{#2}\label{fig:#1}
	\end{figure}
}



% El comando \figura nos permite insertar figuras comodamente, y utilizando
% siempre el mismo formato. Los parametros son:
% 1 -> Porcentaje del ancho de página que ocupará la figura (de 0 a 1)
% 2 --> Fichero de la imagen
% 3 --> Texto a pie de imagen
% 4 --> Etiqueta (label) para referencias
% 5 --> Opciones que queramos pasarle al \includegraphics
% 6 --> Opciones de posicionamiento a pasarle a \begin{figure}
\newcommand{\figuraConPosicion}[6]{%
  \setlength{\anchoFloat}{#1\textwidth}%
  \addtolength{\anchoFloat}{-4\fboxsep}%
  \setlength{\anchoFigura}{\anchoFloat}%
  \begin{figure}[#6]
    \begin{center}%
      \Ovalbox{%
        \begin{minipage}{\anchoFloat}%
          \begin{center}%
            \includegraphics[width=\anchoFigura,#5]{#2}%
            \caption{#3}%
            \label{#4}%
          \end{center}%
        \end{minipage}
      }%
    \end{center}%
  \end{figure}%
}

%
% Comando para incluir imágenes en formato apaisado (sin marco).
\newcommand{\figuraApaisadaSinMarco}[5]{%
  \begin{figure}%
    \begin{center}%
    \includegraphics[angle=90,height=#1\textheight,#5]{#2}%
    \caption{#3}%
    \label{#4}%
    \end{center}%
  \end{figure}%
}
% Para las tablas
\newcommand{\otoprule}{\midrule [\heavyrulewidth]}
%
% Nuevo comando para tablas pequeñas (menos de una página).
\newcommand{\tablaSmall}[5]{%
 \begin{table}
  \begin{center}
   \rowcolors {2}{gray!35}{}
   \begin{tabular}{#2}
    \toprule
    #4
    \otoprule
    #5
    \bottomrule
   \end{tabular}
   \caption{#1}
   \label{tabla:#3}
  \end{center}
 \end{table}
}

%
%Para el float H de tablaSmallSinColores
\usepackage{float}

%
% Nuevo comando para tablas pequeñas (menos de una página).
\newcommand{\tablaSmallSinColores}[5]{%
 \begin{table}[H]
  \begin{center}
   \begin{tabular}{#2}
    \toprule
    #4
    \otoprule
    #5
    \bottomrule
   \end{tabular}
   \caption{#1}
   \label{tabla:#3}
  \end{center}
 \end{table}
}

\newcommand{\tablaApaisadaSmall}[5]{%
\begin{landscape}
  \begin{table}
   \begin{center}
    \rowcolors {2}{gray!35}{}
    \begin{tabular}{#2}
     \toprule
     #4
     \otoprule
     #5
     \bottomrule
    \end{tabular}
    \caption{#1}
    \label{tabla:#3}
   \end{center}
  \end{table}
\end{landscape}
}

%
% Nuevo comando para tablas grandes con cabecera y filas alternas coloreadas en gris.
\newcommand{\tabla}[6]{%
  \begin{center}
    \tablefirsthead{
      \toprule
      #5
      \otoprule
    }
    \tablehead{
      \multicolumn{#3}{l}{\small\sl continúa desde la página anterior}\\
      \toprule
      #5
      \otoprule
    }
    \tabletail{
      \hline
      \multicolumn{#3}{r}{\small\sl continúa en la página siguiente}\\
    }
    \tablelasttail{
      \hline
    }
    \bottomcaption{#1}
    \rowcolors {2}{gray!35}{}
    \begin{xtabular}{#2}
      #6
      \bottomrule
    \end{xtabular}
    \label{tabla:#4}
  \end{center}
}

%
% Nuevo comando para tablas grandes con cabecera.
\newcommand{\tablaSinColores}[6]{%
  \begin{center}
    \tablefirsthead{
      \toprule
      #5
      \otoprule
    }
    \tablehead{
      \multicolumn{#3}{l}{\small\sl continúa desde la página anterior}\\
      \toprule
      #5
      \otoprule
    }
    \tabletail{
      \hline
      \multicolumn{#3}{r}{\small\sl continúa en la página siguiente}\\
    }
    \tablelasttail{
      \hline
    }
    \bottomcaption{#1}
    \begin{xtabular}{#2}
      #6
      \bottomrule
    \end{xtabular}
    \label{tabla:#4}
  \end{center}
}

%
% Nuevo comando para tablas grandes sin cabecera.
\newcommand{\tablaSinCabecera}[5]{%
  \begin{center}
    \tablefirsthead{
      \toprule
    }
    \tablehead{
      \multicolumn{#3}{l}{\small\sl continúa desde la página anterior}\\
      \hline
    }
    \tabletail{
      \hline
      \multicolumn{#3}{r}{\small\sl continúa en la página siguiente}\\
    }
    \tablelasttail{
      \hline
    }
    \bottomcaption{#1}
  \begin{xtabular}{#2}
    #5
   \bottomrule
  \end{xtabular}
  \label{tabla:#4}
  \end{center}
}



\definecolor{cgoLight}{HTML}{EEEEEE}
\definecolor{cgoExtralight}{HTML}{FFFFFF}

%
% Nuevo comando para tablas grandes sin cabecera.
\newcommand{\tablaSinCabeceraConBandas}[5]{%
  \begin{center}
    \tablefirsthead{
      \toprule
    }
    \tablehead{
      \multicolumn{#3}{l}{\small\sl continúa desde la página anterior}\\
      \hline
    }
    \tabletail{
      \hline
      \multicolumn{#3}{r}{\small\sl continúa en la página siguiente}\\
    }
    \tablelasttail{
      \hline
    }
    \bottomcaption{#1}
    \rowcolors[]{1}{cgoExtralight}{cgoLight}

  \begin{xtabular}{#2}
    #5
   \bottomrule
  \end{xtabular}
  \label{tabla:#4}
  \end{center}
}




\graphicspath{ {./img/} }

% Capítulos
\chapterstyle{bianchi}
\newcommand{\capitulo}[2]{
	\setcounter{chapter}{#1}
	\setcounter{section}{0}
	\chapter*{#2}
	\addcontentsline{toc}{chapter}{#2}
	\markboth{#2}{#2}
}

% Apéndices
\renewcommand{\appendixname}{Apéndice}
\renewcommand*\cftappendixname{\appendixname}

\newcommand{\apendice}[1]{
	%\renewcommand{\thechapter}{A}
	\chapter{#1}
}

\renewcommand*\cftappendixname{\appendixname\ }

% Formato de portada
\makeatletter
\usepackage{xcolor}
\newcommand{\tutor}[1]{\def\@tutor{#1}}
\newcommand{\course}[1]{\def\@course{#1}}
\definecolor{cpardoBox}{HTML}{E6E6FF}
\def\maketitle{
  \null
  \thispagestyle{empty}
  % Cabecera ----------------
\noindent\includegraphics[width=\textwidth]{cabecera}\vspace{1cm}%
  \vfill
  % Título proyecto y escudo informática ----------------
  \colorbox{cpardoBox}{%
    \begin{minipage}{.8\textwidth}
      \vspace{.5cm}\Large
      \begin{center}
      \textbf{TFG del Grado en Ingeniería Informática}\vspace{.6cm}\\
      \textbf{\LARGE\@title{}}
      \end{center}
      \vspace{.2cm}
    \end{minipage}

  }%
  \hfill\begin{minipage}{.20\textwidth}
    \includegraphics[width=\textwidth]{escudoInfor}
  \end{minipage}
  \vfill
  % Datos de alumno, curso y tutores ------------------
  \begin{center}%
  {%
    \noindent\LARGE
    Presentado por \@author{}\\ 
    en Universidad de Burgos --- \@date{}\\
    Tutor: \@tutor{}\\
  }%
  \end{center}%
  \null
  \cleardoublepage
  }
\makeatother


% Datos de portada
\title{título del TFG \\Documentación Técnica}
\author{nombre alumno}
\tutor{nombre tutor}
\date{\today}

\begin{document}

\maketitle



\cleardoublepage



%%%%%%%%%%%%%%%%%%%%%%%%%%%%%%%%%%%%%%%%%%%%%%%%%%%%%%%%%%%%%%%%%%%%%%%%%%%%%%%%%%%%%%%%



\frontmatter


\clearpage

% Indices
\tableofcontents

\clearpage

\listoffigures

\clearpage

\listoftables

\clearpage

\mainmatter

\appendix

\apendice{Plan de Proyecto Software}

\section{Introducción}

Antes de profundizar en el desarrollo del proyecto es importante mencionar las circunstancias exteriores que rodean a este trabajo. El trabajo se comenzó a desarrollar en Marzo, ya avanzado el curso, por motivos de trabajo. A esto se le añade que al poco de comenzar, el COVID-19 llegó a nuestro país y puso todo patas arriba, haciendo que las comunicaciones con las partes interesadas en la aplicación fueran muy complicadas, hasta que todo se empezó a estabilizar un poco. 


\section{Planificación temporal}

Contando con todos lo problemas que se tuvieron desde el principio, se optó por utilizar una metodología ágil basada en la programación extrema. La programación extrema se diferencia de otras principalmente en que pone más énfasis en la adaptabilidad que en la previsibilidad, así que teniendo en cuenta todos los factores que nos rodeaban, optamos por esta.

Semanalmente se realizaban entre una y dos reuniones con el tutor, para establecer requisitos y objetivos (sobre todo al principio) y resolver posibles dudas. En estas reuniones se establecía un objetivo basado en las circunstancias generalmente, para la semana siguiente. 

Cada semana se realizaba un commit con los cambios realizados, y se actualizaba un archivo de documentación temporal, en el que quedan anotados todos los pasos que se han ido realizando.

No obstante, todo el ciclo del proyecto se puede dividir en diferentes etapas que se definirán a continuación.

\subsection{Planificación y formación}

La primera parte del proceso que lleno prácticamente un mes de trabajo, consistió en la planificación inicial del proyecto. Se investigaron las posibles herramientas que se podrían utilizar, y se decidió cuáles eran las que se iban a utilizar. La aplicación se desarrollaría en lenguaje java con Eclipse, usando JSP y utilizando Maven para la configuración. Esta elección se tomó en base a los conocimientos que ya se tenían sobre el programa y sobre el lenguaje, adquiridos a lo largo de la carrera. La base de datos se pasaría de SQL Server a MySQL ya que tiene una licencia de código abierto y tiene una integración sencilla y funcional con Java y Eclipse.

Luego se comenzó un proceso de autoformación, en el para empezar que se estudió todo el comportamiento de la aplicación ya existente, así como luego se estudió el funcionamiento de todas las herramientas que se iban a utilizar, realizando pruebas y estudiando sus documentaciones.

Por los motivos ya mencionados previamente, esta fase se alargo más de lo que se esperaba, quedando a la espera de los materiales necesarios, por lo que se dedico aprovechar el tiempo con videotutoriales y cursos.


Se realizaron pruebas en Eclipse usando Maven, y siguiendo un tutorial en Udemy para refrescar los conocimientos \cite{curso:JSPin25Steps}.

Para conectar Java con MySQL, se encontró un videotutorial que solucionó algunas dudas sobre su configuración \cite{tutorial:javaMySQL}


\subsection{Importación y adaptación}

Cuando finalmente recibimos la información del servidor, lo recibimos en un formato de imagen de disco proveniente de Ghost Norton en formato VHD. Esto nos dejaba con una copia del disco duro en el que estaban las tablas que necesitábamos.
Para montar la imagen utilizamos Virtual Vox NT4, pero por alguna razón la maquina no arrancaba. Tras muchas pruebas encontramos que la causa era la tarjeta de red, no era compatible. Tras esto se realizó una búsqueda de una compatible. Finalmente se utilizó AMD PC net family. 

Con todo listo se pudo ejecutar la imagen, y a través de admin SQL exportamos los datos que necesitábamos.

Una vez obtenidas las tablas extraidas, se realizó la instalación de MYSQL descagando el software de su página oficial. Desde la página buscaremos el MySQL Comunnity Server y después seleccionar la plataforma en la que lo vamos a instalar.

Una vez tengamos el ejecutable solo hay que abrirlo y seleccionar el modo de instalación. En mi caso se seleccionó Developer Default ya que se quería instalar tanto MySQL Server como MySql Workbech, un software que proporciona diversas opciones de administraciónde de bases de datos relacionales así como herramientas que nos permiten utilizar consultas SQL mediante un entorno visual.

\subsection{Requisitos funcionales}

La tarea de recoger unos nuevos requisitos funcionales claros fue complicada. Con todo lo sucedido por el confinamiento, las reuniones entre los usuarios de la aplicación fueron imposibles. debido a ello, la única manera de comunicarse fue a través de una serie de correos que los propios usuarios enviaron a Gonzalo. Pedro Renedo recopiló esos correos y me los envío a mi.

En los correos la información era poco específica y caótica, por lo que hizo falta realizar una tarea de análisis para poder extraer unos requisitos funcionales claros. 

Tras esta tarea de análisis e investigación tanto de la aplicación como de conceptos de química utilizados, se extrajeron unas historias de usuario a medio completar, con el objetivo de que o bien Gonzalo, o bien los usuarios, aportaran más información, y sobre todo información más específica. 

Al no recibir respuesta por parte de la universidad (de nuevo por razones del confinamiento) se decidió investigar los requisitos por cuenta propia y establecer unos requisitos temporales para poder ir trabajando sobre ellos. Los requisitos finalmente elegidos no eran definitivos y quedaban abiertos a posibles cambios en el futuro.


\subsection{Diseño de la aplicación}

Para realizar el diseño de la aplicación se comenzó por crear un logotipo, elegir unos colores y una tipografía. Con todo ello, se creo un primer prototipo de la página de inicio (login). Para la creación tanto del logotipo como del primer diseño de la página de login se ha utilizado la herramienta Illustrator. 

Con un diseño gráfico ya establecido, se creó inicialmente la página de login utilizando bootstrap. Inicialmente solo se implementó la parte gráfica de la aplicación, sirviendo a modo de práctica en la utilización de bootstrap. Para ello se realizó una formación en bootstrap a través de la documentación de su página web, y de diversos tutoriales de los que disponen \cite{doc:Bootstrap4}.

Se descargaron todos los archivos pertinentes (css's y javascripts) y se importaron en el proyecto. Para trasladar el diseño creado en illustrator, se ha utilizado una herramienta intermedia llamada Figma.

Con el proyecto ya más avanzado, se creó todo el diseño gráfico del resto de ventanas a utilizar. También se refactorizó el código CSS para ajustarse a una metodología BEM, con el fin de mejorar su legibilidad y reutilización.

También se añadió una libreria javascript llamada animation.js. Esta librería es una biblioteca de animación JavaScript que funciona con propiedades CSS, transformaciones CSS individuales, SVG o cualquier atributo DOM y objetos JavaScript.

Con ella se creó una animación muy simple para la página de login a modo de prueba.


\subsection{Desarrollo de la aplicación}

La fase desarrollo de la aplicación fue la que más tiempo llevó, y la que mas actualizaciones en el proyecto supuso. Lo primero que se realizó fue la conexión de Java con MySQL. Se importó en el proyecto el conector java-mysql, que permite realizar consultas en la base de datos desde java y se creo la clase conexión para realizar todas las configuraciones pertinentes a la conexión. Esta misma clase Conexión, es la que termino albergando todos los métodos  que realizan alguna comunicación con la base de datos.

Con ello se creó la página de login, que al ser la primera, sirvió a modo de practica con los lenguajes y herramientas. Después se implemento todo el sistema de mantenimiento de la web, el cual va enfocado solamente al usuario administrador. Se creó un paquete admin para contener todas las clases implicadas en estas tareas. 

Se crearon los objetos Area, Centro, Departamento y Usuario para poder almacenar toda la información que se lee de las respectivas tablas de la base de datos. Además se creo una clase enumeración para el Rol. De esta manera el rol pasa a ser una variable estática definida por la aplicación, ya que no interesa que un usuario administrador pueda crear roles.

Con cada página JSP que se crea, se crea también su correspondiente Servlet, de forma que de primeras se creo el servlet del login y de admin.

Una vez terminadas las funcionalidades de administración y mantenimiento,  se pasó a crear el paquete usuario, que contendría todas las clases pertenecientes a las tablas a las que el usuario puede tener acceso.

Se comenzó por la que denominamos página de gestión, es decir, la página que mostrará todas las entradas y salidas de productos, y que será la primera página que verá el usuario al acceder. Con ella se crearon los objetos java de Ficha, Entrada y Salida, con su correspondiente página JSP y su servlet.

Para el funcionamiento de la página de gestión, también era necesaria la implementación del resto de páginas a las que denominamos auxiliares, que son las que permiten al usuario gestionar datos de la base de datos como son: las calidades, las marcas, los productos, los proveedores y las salidas.

Cada tabla de las anteriores dispone de su propia página JSP y se servlet, ya que el usuario podrá tanto añadir como editar los datos de todas ellas.


\subsection{Cambios y arreglos}

Una vez se obtuvieron prácticamente todas las funcionalidades básicas de la aplicación, se dedico el tiempo a arreglar y perfeccionar cosas que iban quedando en la recámara, como la encriptacion de las contraseñas, o el correcto funcionamiento de los filtros que poseen las páginas. 

Además se realizó una reunión con los usuarios del laboratorio que van a utilizar la aplicación, y se les mostró una preview de lo que iba funcionando hasta el momento. Tras esta reunión se plantearon unos nuevos puntos de vista, que requerían de algún cambio en el funcionamiento del sistema que, se añadieron y pasaron a implementarse.

El cambio más importante a destacar tras esta reunión, es el cambió drástico de la página principal, que pasó de mostrar las entradas por un lado, y las salidas por otro, a mostrar directamente las fichas de producto completas junto con sus existencias actuales.

Este cambio produjo también un cambio importante en la base de datos y en el funcionamiento en general, ya que se paso de tratar a las entradas y salidas de una manera primordial, a centrarse en las fichas y sus existencias. 

Este cambio que fue bastante necesario, ya que se llevaba tiempo planteando, pero al no poder resolver las dudas con el equipo, no se había realizado aún, hasta que finalmente en la reunión se resolvieron las dudas.

\section{Estudio de viabilidad}

En el siguiente apartado se analizarán todos los parámetros que pueden estar involucrados a la hora de estudiar la viabilidad del proyecto que nos ocupa.

\subsection{Viabilidad económica}

En el siguiente apartado se realizará un estudio de la viabilidad económica del proyecto, suponiendo que el proyecto se ha realizado en el mundo laboral y no en el plan de estudios del grado.

\subsubsection{Coste de personal}
Por un lado hay que contar los recursos humanos empleados en el desarrollo de nuestro proyecto, que serían el alumno desarrollador y el tutor.

El alumno ha empleado un total de 17 semanas a jornada completa (40 horas semanales). Suponiendo que el alumno tiene un sueldo mensual bruto de 1200€ , a lo que hay que restarle los gastos de seguridad social \footnote{\url{http://www.seg-social.es/wps/portal/wss/internet/Trabajadores/CotizacionRecaudacionTrabajadores/36537}}:

\begin{itemize}
	\item Desempleo: 1.55\%
	\item Formación Profesional: 0.10\%
	\item Contingencias Comunes: 4.7\%
\end{itemize}

Por lo que si hacemos la cuenta:

\[1200\euro - 1200\euro \times \left (\frac{1.55 + 0.10 + 4.7}{100}\right ) = 1123.8\euro \]

Así que, multiplicando por el número de meses que ha durado el proyecto 4 x 1123.8\euro hace un total de \textbf{4495.2€} con respecto al alumno.

Además del alumno, hay que tener en cuenta el trabajo realizado por parte del tutor. Para ello contaremos con que una hora de trabajo son 50€ brutos la hora incluyendo instalaciones, costes laborales y otros indirectos como gastos generales, equipamiento, espacios, etc.

Teniendo en cuenta que se han realizado entre una y dos reuniones semanales de 45 minutos, otras tres reuniones con la facultad de 45 minutos, y sumando además unas 10 horas extra para revisar documentación, mediación con el decano y consultas procedentes del desarrollador, nos sale un total de:

\[ \left (  \left ( 17\: sem. \times \frac{1.5\, reu.\times 45\, mins}{60} \right )+\left ( 3\, reu. \times \frac{45\, mins}{60} \right ) + 10\right )\times\, 50\euro/h = 1568.75\euro \]


\subsubsection{Coste de hardware}

Para calcular el gasto de hardware, hay que valorar el material utilizado durante el proyecto por el alumno:
\begin{itemize}
	\item Ordenador potátil: 900€
	\item Monitor: 100€
	\item Teclado y ratón: 25€
\end{itemize}

En total tenemos un gasto de 1025€. Según la agencia tributaria este tipo de material tiene una vida útil de 10 años \footnote{\url{https://www.agenciatributaria.es/AEAT.internet/Inicio/_Segmentos_/Empresas_y_profesionales/Empresas/Impuesto_sobre_Sociedades/Periodos_impositivos_a_partir_de_1_1_2015/Base_imponible/Amortizacion/Tabla_de_coeficientes_de_amortizacion_lineal_.shtml}}, y nuestro proyecto ha durado un total de 17 semanas por lo que el coste de amortización sería:

\[\frac{1025\euro}{10\, a\textit{ñ}os} \times \frac{1\, a\textit{ñ}o}{52\, semanas} \times 17\, semanas \approx  34\euro \]

\subsubsection{Coste de software}

Para calcular los gastos de software contamos con que practicamente todos los programas utilizados son de licencia gratuita, a excepción de Adobe Illustrator, por lo que hay que contar con un coste por su licencia de 29.99€ al mes (sin iva). Hay que tener en cuenta que la licencia se paga por mes completo, por lo que se realizará el cálculo sobre 20 semanas en vez de 17.

\[ (29.99 \euro + 21\% \, iva) \times 5\, meses = 181.45 \euro\]

A esto habría que añadirle el coste del Office 365, que eligiendo el paquete Microsoft 365 Empresa Estándar nos sale a 10.50€ al mes por usuario (sin iva), pero con una suscripción anual.

\[ ((10.50 \euro + 21\% \, iva) \times 12\, meses) \times 2\, usuarios = 305.04 \euro\]

\subsubsection{Coste total}

Con todos los costes mencionados, podemos ver todos los gastos y su suma total en la tabla~\ref{tabla:costes}

\begin{table}[h]
	\centering
	\caption{Coste total del proyecto}
	\label{tabla:costes}
	\rowcolors {2}{gray!20}{}
	\begin{tabular}{p{6cm} r}
		\toprule
		Costes & Importe \\ \midrule
		Desarrollador			     & 4495.2\euro   						\\
		Tutor  				         & 1568.75\euro							\\
		Hardware				     & 34\euro								\\
		Software 				     & 486.49\euro 							\\
		\textbf{Total} 			     & 6584.44\euro 						\\  
		\bottomrule
	\end{tabular}
\end{table}


\subsection{Viabilidad legal}
 Para la viabilidad legal hay que analizar las librerías que se han utilizado, buscando sus licencias, analizando las compatibilidades y buscando las más restrictiva.
 
En la tabla~\ref{tabla:licencias} se muestra el listado de todas las licencias.

Al ser la licencia \textit{GNU General Public Licence} la más restrictiva, hay que calificar a Geslab 2.0 bajo una licencia \textit{GNU General Public Licence}.

\begin{table}[h]
	\centering
	\caption{Herramientas utilizadas y sus licencias}
	\label{tabla:licencias}
	\rowcolors {2}{gray!20}{}
	\begin{tabular}{p{4cm} p{6cm}}
		\toprule
		Herramienta & Licencia \\ \midrule
		Bootstrap 4				     & MIT   								\\
		Animation.js		         & MIT   								\\
		jQuery				         & MIT   								\\
		Font Awesome			     & MIT 									\\ 
		MySQL Connector				 & GNU General Public Licence			\\ \bottomrule
	\end{tabular}
\end{table}



\apendice{Especificación de Requisitos}

\section{Introducción}

En el siguiente apartado se presentarán los diferentes objetivos a cumplir así como los requisitos funcionales que han dado funcionalidad a nuestra aplicación.

\section{Objetivos generales}

los objetivos principales del proyecto son:

\begin{itemize}
	\item Recrear la aplicación Geslab 1.0 como una aplicación web.
	\item Mantener todas las funcionalidades que ya poseía.
	\item Añadir distintos tipos de usuarios: administrador, gestor y usuario.
	\item Añadir nuevos campos a la base de datos tales como: pureza, peso molecular, N.º CAS, ...
	\item Implementar un sistema de búsqueda intuitivo y funcional
	\item La aplicación utilizará lenguaje Java
	\item La base de datos utilizará MySql
\end{itemize}

\section{Catálogo de requisitos}

Los requisitos que debe de satisfacer la aplicación son los siguientes.

\subsection{Requisitos funcionales}

\begin{itemize}
\item \textbf{RF-1 Accesibilidad:} La aplicación deberá de ser accesible desde cualquier navegador web de la universidad.

\item \textbf{RF-2 Administración:} La aplicación deberá permitir gestionar sus propios datos sin necesidad de alterar a mano la base de datos
	\begin{itemize}
	\item\textbf{ RF-2.1} Deberán de existir tres roles en la aplicación: administrador, gestor de inventario y usuario.
	
	\item \textbf{RF-2.2} El usuario solamente puede ver las existencias, sin poder modificar o añadir.
	
	\item \textbf{RF-2.3} El gestor de inventario puede añadir existencias y editar la información.
	
	\item \textbf{RF-2.4} El administrador puede gestionar los usuarios, las áreas, los departamentos y los centros.
	
	\item \textbf{RF-2.5} Los usuarios tienen que tener asignado un rol, un área a la que pertenecen, si es un usuario con cuenta federada o no, y si la cuenta esta activa.
	
	\item \textbf{RF-2.6} El administrador puede añadir y editar usuarios, pero no eliminar, por seguridad. Para ello está la opción de dejar al usuario inactivo.
	
	\item \textbf{RF-2.7} Cada área debe de gestionar solamente sus propios productos.
	
	\item \textbf{RF-2.8} Cada área puede tener varias ubicaciones donde se pueden guardar los productos (armarios, estanterías...)
	
	\item \textbf{RF-2.9} Cada área puede ocultar sus ubicaciones o no al resto de áreas, de forma que las otras áreas no puedan conocer las existencias de los productos almacenados en esas ubicaciones.
	
	\end{itemize}
	
\item \textbf{RF-3 Gestión de productos}
	\begin{itemize}
	\item\textbf{ RF-3.1} Se almacenarán los productos con su N.º Cas, fórmula, fórmula desarrollada, peso molecular, pureza, N.º EINECS, N.º EC, precauciones y su hoja de seguridad(MSDS). De estos campos solamente serán obligatorios el nº cas y la fórmula.
	
	\item\textbf{ RF-3.2} Los productos serán comunes a todas las áreas.
	
	\item\textbf{ RF-3.3} Cada producto puede tener varios nombres (distintas nomenclaturas, en inglés, en español), por lo que se deberán almacenar también.
	
	\item\textbf{ RF-3.4} Al tener varios nombres, cada producto se diferenciará por su N.º CAS, al ser este único.
	
	\item\textbf{ RF-3.5} Cada producto además almacenará una serie de peligros, prudencias y pictogramas de seguridad asociados al mismo. Estos peligros prudencias y pictogramas han de ser comunes a todos los productos.
	
	\end{itemize}
	
\item \textbf{RF-4 Gestión de fichas}
	\begin{itemize}
	\item\textbf{ RF-4.1} Las existencias de los productos se gestiona a través de fichas. Cada ficha corresponde a un producto, de una capacidad (en gramos o mililitros), guardado en una 	ubicación, de una marca y un proveedor, con una calidad determinada, un N.º de lote y una caducidad. Además, se podrá indicar si cada ficha corresponde a un residuo o no. 
	
	\item\textbf{ RF-4.2} Todos los atributos de la ficha son obligatorios, ya que es la manera de diferenciarlos (no es lo mismo un bote de Helio que caduca en 2025 que un bote de Helio que caduca en 2022, por ejemplo).

	\item\textbf{ RF-4.3} De cada ficha puede haber entradas o salidas, las cuales se almacenarán con la ficha a la que pertenecen, las unidades que entran o salen, la fecha, una nota y el usuario que la realiza.
	
	\item\textbf{ RF-4.4} Cada ficha tiene un stock, que es actualizado cada vez que se realiza una entrada o una salida.
	
	\item\textbf{ RF-4.5} Los usuarios pueden consultar el stock de todas las fichas de otras áreas siempre y cuando no se encuentren en una ubicación oculta.
	
	\item\textbf{ RF-4.6} Se podrá realizar un filtrado de las fichas para poder realizar una búsqueda más rápida.
	
	\item\textbf{ RF-4.7} Cada área podrá gestionar sus propias fichas, pudiendo añadir fichas o añadir entradas o salidas. No se permite el eliminado de fichas, entradas o salidas por seguridad.
	\end{itemize} 
	
	\item \textbf{RF-5 Gestión de marcas y proveedores}
	\begin{itemize}
	\item\textbf{ RF-5.1} Se almacenarán las marcas y los proveedores, junto con información como el nombre, la dirección, el teléfono, el fax y el email. Las marcas y los proveedores son comunes a todos los usuarios.
	
	\item\textbf{ RF-5.2} También se almacenará la relación entre estos. Un proveedor puede llevar varias marcas, y una marca puede ser llevada por varios proveedores.
	
	\item\textbf{ RF-5.3} Cualquier área puede añadir o editar los proveedores y las marcas, así como la relación entre estas.
	\end{itemize}
	
	\item \textbf{RF-6 Gestión de calidades}
	\begin{itemize}
	\item\textbf{ RF-6.1} También se almacenarán las calidades, siendo comunes a todos los usuarios y áreas.
	
	\item\textbf{ RF-6.2} Cualquier área puede añadir o editar las calidades.
	
	\end{itemize}
	
	\item \textbf{RF-7 Gestión de ubicaciones}
	\begin{itemize}
	\item\textbf{ RF-7.1} Se almacenarán las ubicaciones junto con su nombre, su área, su centro y si es oculta o no.
	
	\item\textbf{ RF-7.2} Cada área podrá gestionar sus propias ubicaciones, pero no las de las demás áreas. Las ubicaciones se pueden añadir o editar.
	
	\end{itemize}
	
\end{itemize}


\subsection{Requisitos no funcionales}

\begin{itemize}
\item \textbf{RNF-1 Usabilidad:} La aplicación debe de ser intuitiva y cómoda de utilizar, usando elementos visuales conocidos para el usuario medio

\item \textbf{RNF-2 Rendimiento:} El rendimiento de la aplicación ha de ser óptimo, haciendo que su navegación sea fluida.

\item \textbf{RNF-3 Mantenibilidad:} La aplicación debe de facilitar el mantenimiento y la posible incorporación de nuevas características.

\item \textbf{RNF-4 Despliegue:} La aplicación debe de estar desplegada de forma que pueda ser accesible sin necesidad de realizar una ejecución local.

\end{itemize}


\section{Especificación de requisitos}

\subsection{Actores}

En nuestro caso existen cuatro actores, el administrador, el gestor de inventario, el usuario y la base de datos.

\subsection{Diagramas de casos de uso}

Veamos a continuación el diagrama de casos de uso de nuestro proyecto por niveles en las siguientes figuras.

\imagen{Nivel-1}{Nivel 1 del diagrama de casos de uso.} {1.2}
\imagen{Nivel-1A}{Nivel 1A del diagrama de casos de uso.} {1.1}
\imagen{Nivel-1B}{Nivel 1B del diagrama de casos de uso.} {1.2}
\imagen{Nivel-1C}{Nivel 1C del diagrama de casos de uso.} {0.9}

%Gestion de centros -----------------------------
\begin{table}[h]
	\centering
	\label{tabla:cu1}
	\begin{tabular}{@{}
		>{\columncolor[HTML]{FFFFFF}}p {.25\textwidth} p {.75\textwidth}@{}}
		\toprule
		\textbf{Caso de uso 1}   & \textbf{Visualizar centros} \\ \midrule
		\textbf{Versión}     & 2.0 \\ \midrule
		\textbf{Requisitos}	&  RF-2.4 \\ \midrule
		\textbf{Descripción}     & Permite al administrador visualizar los centros existentes. \\ \midrule
		\textbf{Precondiciones}  & 
		\begin{compactitem}
			\item El usuario se ha logueado en la aplicación con rol de administrador. 
		\end{compactitem}
		 \\ \midrule
		\textbf{Acciones} & 
		El administrador selecciona la pestaña Centros. 
		\\ \midrule
		\textbf{Postcondiciones} & -  \\ \midrule
		\textbf{Excepciones} &   - \\ \midrule
		\textbf{Importancia}     & Media \\ \bottomrule
	\end{tabular}
	\caption{Caso de uso 1 - Visualizar centros}
\end{table}

\begin{table}[h]
	\centering
	\label{tabla:cu2}
	\begin{tabular}{@{}
		>{\columncolor[HTML]{FFFFFF}}p {.25\textwidth} p {.75\textwidth}@{}}
		\toprule
		\textbf{Caso de uso 2}   & \textbf{Editar centros} \\ \midrule
		\textbf{Versión}     & 2.0 \\ \midrule
		\textbf{Requisitos}	&  RF-2.4 \\ \midrule
		\textbf{Descripción}     & Permite al administrador editar los centros existentes. \\ \midrule
		\textbf{Precondiciones}  & 
		\begin{compactitem}
			\item El usuario se ha logueado en la aplicación con rol de administrador. 
		\end{compactitem}
		 \\ \midrule
		\textbf{Acciones} & 
		El usuario tiene que clickar sobre el icono de la derecha de la fila del centro que quiera editar. Después escribirá el nombre que quiere poner y pulsará en el botón de confirmar que aparece al lado del de editar. 
		\\ \midrule
		\textbf{Postcondiciones} & La aplicación volverá a mostrar la lista con el centro  editado. \\ \midrule
		\textbf{Excepciones} & Si el nombre que inserta ya es igual al de un centro existente, el sistema mostrará un mensaje de error. \\ \midrule
		\textbf{Importancia}     & Media \\ \bottomrule
	\end{tabular}
	\caption{Caso de uso 2 - Editar centros}
\end{table}

\begin{table}[h]
	\centering
	\label{tabla:cu3}
	\begin{tabular}{@{}
		>{\columncolor[HTML]{FFFFFF}}p {.25\textwidth} p {.75\textwidth}@{}}
		\toprule
		\textbf{Caso de uso 3}   & \textbf{Insertar centros} \\ \midrule
		\textbf{Versión}     & 2.0 \\ \midrule
		\textbf{Requisitos}	&  RF-2.4 \\ \midrule
		\textbf{Descripción}     & Permite añadir un nuevo centro. \\ \midrule
		\textbf{Precondiciones}  & 
		\begin{compactitem}
			\item El usuario se ha logueado en la aplicación con rol de administrador. 
		\end{compactitem}
		 \\ \midrule
		\textbf{Acciones} & 
		El usuario tiene que clickar sobre el botón de Nuevo centro. Después escribirá el nombre que desea ponerle y pulsará en el botón de confirmar que aparece a la derecha de la fila.
		\\ \midrule
		\textbf{Postcondiciones} & La aplicación volverá a mostrar la lista con el centro  añadido. \\ \midrule
		\textbf{Excepciones} & Si el nombre que inserta ya es igual al de un centro existente, el sistema mostrará un mensaje de error. \\ \midrule
		\textbf{Importancia}     & Media \\ \bottomrule
	\end{tabular}
	\caption{Caso de uso 3 - Insertar centros}
\end{table}

%Gestion de departamentos -----------------------------
\begin{table}[h]
	\centering
	\label{tabla:cu4}
	\begin{tabular}{@{}
		>{\columncolor[HTML]{FFFFFF}}p {.25\textwidth} p {.75\textwidth}@{}}
		\toprule
		\textbf{Caso de uso 4}   & \textbf{Visualizar departamentos} \\ \midrule
		\textbf{Versión}     & 2.0 \\ \midrule
		\textbf{Requisitos}	&  RF-2.4 \\ \midrule
		\textbf{Descripción}     & Permite al administrador visualizar los departamentos existentes. \\ \midrule
		\textbf{Precondiciones}  & 
		\begin{compactitem}
			\item El usuario se ha logueado en la aplicación con rol de administrador. 
		\end{compactitem}
		 \\ \midrule
		\textbf{Acciones} & 
		El administrador selecciona la pestaña Departamentos. 
		\\ \midrule
		\textbf{Postcondiciones} & -  \\ \midrule
		\textbf{Excepciones} &   - \\ \midrule
		\textbf{Importancia}     & Media \\ \bottomrule
	\end{tabular}
	\caption{Caso de uso 4 - Visualizar departamentos}
\end{table}

\begin{table}[h]
	\centering
	\label{tabla:cu5}
	\begin{tabular}{@{}
		>{\columncolor[HTML]{FFFFFF}}p {.25\textwidth} p {.75\textwidth}@{}}
		\toprule
		\textbf{Caso de uso 5}   & \textbf{Editar departamentos} \\ \midrule
		\textbf{Versión}     & 2.0 \\ \midrule
		\textbf{Requisitos}	&  RF-2.4 \\ \midrule
		\textbf{Descripción}     & Permite al administrador editar los departamentos existentes. \\ \midrule
		\textbf{Precondiciones}  & 
		\begin{compactitem}
			\item El usuario se ha logueado en la aplicación con rol de administrador. 
		\end{compactitem}
		 \\ \midrule
		\textbf{Acciones} & 
		El usuario tiene que clickar sobre el icono de la derecha de la fila del departamento que quiera editar. Después escribirá el nombre que quiere poner y pulsará en el botón de confirmar que aparece al lado del de editar. 
		\\ \midrule
		\textbf{Postcondiciones} & La aplicación volverá a mostrar la lista con el departamento  editado. \\ \midrule
		\textbf{Excepciones} & Si el nombre que inserta ya es igual al de un departamento existente, el sistema mostrará un mensaje de error. \\ \midrule
		\textbf{Importancia}     & Media \\ \bottomrule
	\end{tabular}
	\caption{Caso de uso 5 - Editar departamentos}
\end{table}

\begin{table}[h]
	\centering
	\label{tabla:cu6}
	\begin{tabular}{@{}
		>{\columncolor[HTML]{FFFFFF}}p {.25\textwidth} p {.75\textwidth}@{}}
		\toprule
		\textbf{Caso de uso 6}   & \textbf{Insertar departamentos} \\ \midrule
		\textbf{Versión}     & 2.0 \\ \midrule
		\textbf{Requisitos}	&  RF-2.4 \\ \midrule
		\textbf{Descripción}     & Permite añadir un nuevo departamento. \\ \midrule
		\textbf{Precondiciones}  & 
		\begin{compactitem}
			\item El usuario se ha logueado en la aplicación con rol de administrador. 
		\end{compactitem}
		 \\ \midrule
		\textbf{Acciones} & 
		El usuario tiene que clickar sobre el botón de Nuevo departamento. Después escribirá el nombre que desea ponerle y pulsará en el botón de confirmar que aparece a la derecha de la fila.
		\\ \midrule
		\textbf{Postcondiciones} & La aplicación volverá a mostrar la lista con el departamento  añadido. \\ \midrule
		\textbf{Excepciones} & Si el nombre que inserta ya es igual al de un departamento existente, el sistema mostrará un mensaje de error. \\ \midrule
		\textbf{Importancia}     & Media \\ \bottomrule
	\end{tabular}
	\caption{Caso de uso 6 - Insertar departamentos}
\end{table}


%Gestion de áreas -----------------------------
\begin{table}[h]
	\centering
	\label{tabla:cu7}
	\begin{tabular}{@{}
		>{\columncolor[HTML]{FFFFFF}}p {.25\textwidth} p {.75\textwidth}@{}}
		\toprule
		\textbf{Caso de uso 7}   & \textbf{Visualizar áreas} \\ \midrule
		\textbf{Versión}     & 2.0 \\ \midrule
		\textbf{Requisitos}	&  RF-2.4 \\ \midrule
		\textbf{Descripción}     & Permite al administrador visualizar las áreas existentes. \\ \midrule
		\textbf{Precondiciones}  & 
		\begin{compactitem}
			\item El usuario se ha logueado en la aplicación con rol de administrador. 
		\end{compactitem}
		 \\ \midrule
		\textbf{Acciones} & 
		El administrador selecciona la pestaña Áreas. 
		\\ \midrule
		\textbf{Postcondiciones} & -  \\ \midrule
		\textbf{Excepciones} &   - \\ \midrule
		\textbf{Importancia}     & Media \\ \bottomrule
	\end{tabular}
	\caption{Caso de uso 7 - Visualizar áreas}
\end{table}

\begin{table}[h]
	\centering
	\label{tabla:cu8}
	\begin{tabular}{@{}
		>{\columncolor[HTML]{FFFFFF}}p {.25\textwidth} p {.75\textwidth}@{}}
		\toprule
		\textbf{Caso de uso 8}   & \textbf{Editar áreas} \\ \midrule
		\textbf{Versión}     & 2.0 \\ \midrule
		\textbf{Requisitos}	&  RF-2.4 \\ \midrule
		\textbf{Descripción}     & Permite al administrador editar las áreas existentes. \\ \midrule
		\textbf{Precondiciones}  & 
		\begin{compactitem}
			\item El usuario se ha logueado en la aplicación con rol de administrador. 
		\end{compactitem}
		 \\ \midrule
		\textbf{Acciones} & 
		El usuario tiene que clickar sobre el icono de la derecha de la fila del área que quiera editar. Después escribirá el nombre que quiere poner y su departamento y pulsará en el botón de confirmar que aparece al lado del de editar. 
		\\ \midrule
		\textbf{Postcondiciones} & La aplicación volverá a mostrar la lista con el área  editada. \\ \midrule
		\textbf{Excepciones} & Si el nombre que inserta ya es igual al de un área existente, el sistema mostrará un mensaje de error. \\ \midrule
		\textbf{Importancia}     & Media \\ \bottomrule
	\end{tabular}
	\caption{Caso de uso 8 - Editar áreas}
\end{table}

\begin{table}[h]
	\centering
	\label{tabla:cu9}
	\begin{tabular}{@{}
		>{\columncolor[HTML]{FFFFFF}}p {.25\textwidth} p {.75\textwidth}@{}}
		\toprule
		\textbf{Caso de uso 9}   & \textbf{Insertar áreas} \\ \midrule
		\textbf{Versión}     & 2.0 \\ \midrule
		\textbf{Requisitos}	&  RF-2.4 \\ \midrule
		\textbf{Descripción}     & Permite añadir un nuevo área. \\ \midrule
		\textbf{Precondiciones}  & 
		\begin{compactitem}
			\item El usuario se ha logueado en la aplicación con rol de administrador. 
		\end{compactitem}
		 \\ \midrule
		\textbf{Acciones} & 
		El usuario tiene que clickar sobre el botón de Nuevo Área. Después escribirá el nombre que desea ponerle y su departamento y pulsará en el botón de confirmar que aparece a la derecha de la fila.
		\\ \midrule
		\textbf{Postcondiciones} & La aplicación volverá a mostrar la lista con el área añadida. \\ \midrule
		\textbf{Excepciones} & Si el nombre que inserta ya es igual al de un departamento existente, el sistema mostrará un mensaje de error. \\ \midrule
		\textbf{Importancia}     & Media \\ \bottomrule
	\end{tabular}
	\caption{Caso de uso 9 - Insertar áreas}
\end{table}

%Gestion de usuarios -----------------------------
\begin{table}[h]
	\centering
	\label{tabla:cu10}
	\begin{tabular}{@{}
		>{\columncolor[HTML]{FFFFFF}}p {.25\textwidth} p {.75\textwidth}@{}}
		\toprule
		\textbf{Caso de uso 10}   & \textbf{Visualizar usuarios} \\ \midrule
		\textbf{Versión}     & 2.0 \\ \midrule
		\textbf{Requisitos}	&  RF-2.4, RF-2.5\\ \midrule
		\textbf{Descripción}     & Permite al administrador visualizar los usuarios existentes. \\ \midrule
		\textbf{Precondiciones}  & 
		\begin{compactitem}
			\item El usuario se ha logueado en la aplicación con rol de administrador. 
		\end{compactitem}
		 \\ \midrule
		\textbf{Acciones} & 
		El administrador selecciona la pestaña Usuarios. 
		\\ \midrule
		\textbf{Postcondiciones} & -  \\ \midrule
		\textbf{Excepciones} &   - \\ \midrule
		\textbf{Importancia}     & Media \\ \bottomrule
	\end{tabular}
	\caption{Caso de uso 10 - Visualizar usuarios}
\end{table}

\begin{table}[h]
	\centering
	\label{tabla:cu11}
	\begin{tabular}{@{}
		>{\columncolor[HTML]{FFFFFF}}p {.25\textwidth} p {.75\textwidth}@{}}
		\toprule
		\textbf{Caso de uso 11}   & \textbf{Editar usuarios} \\ \midrule
		\textbf{Versión}     & 2.0 \\ \midrule
		\textbf{Requisitos}	&  RF-2.4, RF-2.6 \\ \midrule
		\textbf{Descripción}     & Permite al administrador editar los usuarios existentes. \\ \midrule
		\textbf{Precondiciones}  & 
		\begin{compactitem}
			\item El usuario se ha logueado en la aplicación con rol de administrador. 
		\end{compactitem}
		 \\ \midrule
		\textbf{Acciones} & 
		El usuario tiene que clickar sobre el icono de la derecha de la fila del usuario que quiera editar. Después escribirá el nombre que quiere poner, el rol, el área, si es federada y si está activo y pulsará en el botón de confirmar que aparece al lado del de editar. 
		\\ \midrule
		\textbf{Postcondiciones} & La aplicación volverá a mostrar la lista con el usuario  editado. \\ \midrule
		\textbf{Excepciones} & Si el nombre que inserta ya es igual al de un usuario existente o algún dato es inválido, el sistema mostrará un mensaje de error. \\ \midrule
		\textbf{Importancia}     & Media \\ \bottomrule
	\end{tabular}
	\caption{Caso de uso 11 - Editar usuarios}
\end{table}

\begin{table}[h]
	\centering
	\label{tabla:cu12}
	\begin{tabular}{@{}
		>{\columncolor[HTML]{FFFFFF}}p {.25\textwidth} p {.75\textwidth}@{}}
		\toprule
		\textbf{Caso de uso 12}   & \textbf{Insertar usuarios} \\ \midrule
		\textbf{Versión}     & 2.0 \\ \midrule
		\textbf{Requisitos}	&  RF-2.4 \\ \midrule
		\textbf{Descripción}     & Permite añadir un nuevo usuario. \\ \midrule
		\textbf{Precondiciones}  & 
		\begin{compactitem}
			\item El usuario se ha logueado en la aplicación con rol de administrador. 
		\end{compactitem}
		 \\ \midrule
		\textbf{Acciones} & 
		El usuario tiene que clickar sobre el botón de Nuevo usuario. Después escribirá el nombre que desea ponerle, el rol, el área, si es federada y si está activo y pulsará en el botón de confirmar que aparece a la derecha de la fila.
		\\ \midrule
		\textbf{Postcondiciones} & La aplicación volverá a mostrar la lista con el usuario  añadido. \\ \midrule
		\textbf{Excepciones} & Si el nombre que inserta ya es igual al de un usuario existente o algún dato es inválido, el sistema mostrará un mensaje de error. \\ \midrule
		\textbf{Importancia}     & Media \\ \bottomrule
	\end{tabular}
	\caption{Caso de uso 12 - Insertar usuarios}
\end{table}


%Gestion de existencias -----------------------------
\begin{table}[h]
	\centering
	\label{tabla:cu13}
	\begin{tabular}{@{}
		>{\columncolor[HTML]{FFFFFF}}p {.25\textwidth} p {.75\textwidth}@{}}
		\toprule
		\textbf{Caso de uso 13}   & \textbf{Visualizar existencias} \\ \midrule
		\textbf{Versión}     & 2.0 \\ \midrule
		\textbf{Requisitos}	&  RF-4.1, RF-4.5\\ \midrule
		\textbf{Descripción}     & Permite al administrador visualizar las existencias de fichas existentes, siempre y cuando no estén en una ubicación oculta de otra área. \\ \midrule
		\textbf{Precondiciones}  & 
		\begin{compactitem}
			\item El usuario se ha logueado en la aplicación con rol de Gestor de inventario. 
		\end{compactitem}
		 \\ \midrule
		\textbf{Acciones} & 
		\begin{compactitem}
			\item El gestor se dirige a la página de existencias 
		\end{compactitem}
		\\ \midrule
		\textbf{Postcondiciones} & -  \\ \midrule
		\textbf{Excepciones} &   - \\ \midrule
		\textbf{Importancia}     & Alta \\ \bottomrule
	\end{tabular}
	\caption{Caso de uso 13 - Visualizar existencias}
\end{table}

\begin{table}[h]
	\centering
	\label{tabla:cu14}
	\begin{tabular}{@{}
		>{\columncolor[HTML]{FFFFFF}}p {.25\textwidth} p {.75\textwidth}@{}}
		\toprule
		\textbf{Caso de uso 14}   & \textbf{Filtrar existencias} \\ \midrule
		\textbf{Versión}     & 2.0 \\ \midrule
		\textbf{Requisitos}	&  RF-4.6 \\ \midrule
		\textbf{Descripción}     & Permite filtrar las existencias en función de los parámetros: producto, cas, fórmula, departamento, área, centro, ubicación, marca, proveedor, calidad, ubicación oculta o residuo. \\ \midrule
		\textbf{Precondiciones}  & 
		\begin{compactitem}
			\item El usuario se ha logueado en la aplicación con rol de Gestor de inventario.
		\end{compactitem}
		 \\ \midrule
		\textbf{Acciones} & 
		El usuario rellena o selecciona un valor de los desplegables en la columna de filtrado de la izquierda. En cuanto inserte un carácter el filtrado se producirá automáticamente sin pulsar ningún botón. También puede borrar todos los filtros pulsando en el botón Reiniciar filtro.
		\\ \midrule
		\textbf{Postcondiciones} & La aplicación mostrará en la lista únicamente las fichas que coincidan con los parámetros de la búsqueda introducidos. \\ \midrule
		\textbf{Excepciones} & - \\ \midrule
		\textbf{Importancia} & Baja \\ \bottomrule
	\end{tabular}
	\caption{Caso de uso 14 - Filtrar existencias}
\end{table}

\begin{table}[h]
	\centering
	\label{tabla:cu15}
	\begin{tabular}{@{}
		>{\columncolor[HTML]{FFFFFF}}p {.25\textwidth} p {.75\textwidth}@{}}
		\toprule
		\textbf{Caso de uso 15}   & \textbf{Insertar fichas} \\ \midrule
		\textbf{Versión}     & 2.0 \\ \midrule
		\textbf{Requisitos}	&  RF-4.7 \\ \midrule
		\textbf{Descripción}     & Permite añadir una nueva ficha. \\ \midrule
		\textbf{Precondiciones}  & 
		\begin{compactitem}
			\item El usuario se ha logueado en la aplicación con rol de Gestor de inventario. 
		\end{compactitem}
		 \\ \midrule
		\textbf{Acciones} & 
		El usuario tiene que clickar sobre el botón de Nueva ficha. Después elegirá el producto, su capacidad, g o ml, su ubicación, su marca, su proveedor, su calidad, el nº de lote, su caducidad y si es residuo o no y pulsará en el botón añadir del modal.
		\\ \midrule
		\textbf{Postcondiciones} & La aplicación volverá a mostrar la lista con la ficha añadida. \\ \midrule
		\textbf{Excepciones} & Si la ficha que inserta ya existe o algún dato es inválido, el sistema mostrará un mensaje de error. \\ \midrule
		\textbf{Importancia}     & Alta \\ \bottomrule
	\end{tabular}
	\caption{Caso de uso 15 - Insertar fichas}
\end{table}

\begin{table}[h]
	\centering
	\label{tabla:cu16}
	\begin{tabular}{@{}
		>{\columncolor[HTML]{FFFFFF}}p {.25\textwidth} p {.75\textwidth}@{}}
		\toprule
		\textbf{Caso de uso 16}   & \textbf{Añadir entradas} \\ \midrule
		\textbf{Versión}     & 2.0 \\ \midrule
		\textbf{Requisitos}	& RF-4.3, RF-4.7 \\ \midrule
		\textbf{Descripción}     & Permite añadir entradas de una ficha. \\ \midrule
		\textbf{Precondiciones}  & 
		\begin{compactitem}
			\item El usuario se ha logueado en la aplicación con rol de Gestor de inventario.
			\item La ficha debe de estar en una ubicación del mismo área a la que pertenece el usuario
		\end{compactitem}
		 \\ \midrule
		\textbf{Acciones} & 
		El usuario tiene que clickar sobre el icono + en la columna de la derecha de la fila de la ficha que desea. Después elegirá las unidades, la fecha y una nota y pulsará en el botón añadir del modal.
		\\ \midrule
		\textbf{Postcondiciones} & La aplicación volverá a mostrar la lista con el stock de la ficha añadido. \\ \midrule
		\textbf{Excepciones} & Si algún dato es inválido, o deja la fecha vacía el sistema mostrará un mensaje de error. \\ \midrule
		\textbf{Importancia}     & Alta \\ \bottomrule
	\end{tabular}
	\caption{Caso de uso 16 - Añadir entradas}
\end{table}

\begin{table}[h]
	\centering
	\label{tabla:cu17}
	\begin{tabular}{@{}
		>{\columncolor[HTML]{FFFFFF}}p {.25\textwidth} p {.75\textwidth}@{}}
		\toprule
		\textbf{Caso de uso 17}   & \textbf{Añadir salidas} \\ \midrule
		\textbf{Versión}     & 2.0 \\ \midrule
		\textbf{Requisitos}	&  RF-4.3, RF-4.7 \\ \midrule
		\textbf{Descripción}     & Permite añadir salidas de una ficha. \\ \midrule
		\textbf{Precondiciones}  & 
		\begin{compactitem}
			\item El usuario se ha logueado en la aplicación con rol de Gestor de inventario.
			\item La ficha debe de estar en una ubicación de la misma área a la que pertenece el usuario
		\end{compactitem}
		 \\ \midrule
		\textbf{Acciones} & 
		El usuario tiene que clickar sobre el icono - en la columna de la derecha de la fila de la ficha que desea. Después elegirá las unidades, la fecha y una nota y pulsará en el botón añadir del modal.
		\\ \midrule
		\textbf{Postcondiciones} & La aplicación volverá a mostrar la lista con el stock de la ficha restado. \\ \midrule
		\textbf{Excepciones} & Si algún dato es inválido, o deja la fecha vacía el sistema mostrará un mensaje de error. \\ \midrule
		\textbf{Importancia}     & Alta \\ \bottomrule
	\end{tabular}
	\caption{Caso de uso 17 - Añadir salidas}
\end{table}


%Gestion de productos -----------------------------
\begin{table}[h]
	\centering
	\label{tabla:cu18}
	\begin{tabular}{@{}
		>{\columncolor[HTML]{FFFFFF}}p {.25\textwidth} p {.75\textwidth}@{}}
		\toprule
		\textbf{Caso de uso 18}   & \textbf{Visualizar productos} \\ \midrule
		\textbf{Versión}     & 2.0 \\ \midrule
		\textbf{Requisitos}	&  RF-4.1, RF-4.5\\ \midrule
		\textbf{Descripción}     & Permite al administrador visualizar los productos existentes. \\ \midrule
		\textbf{Precondiciones}  & 
		\begin{compactitem}
			\item El usuario se ha logueado en la aplicación con rol de Gestor de inventario. 
		\end{compactitem}
		 \\ \midrule
		\textbf{Acciones} & 
		\begin{compactitem}
			\item El gestor se dirige a la página de productos 
		\end{compactitem}
		\\ \midrule
		\textbf{Postcondiciones} & -  \\ \midrule
		\textbf{Excepciones} &   - \\ \midrule
		\textbf{Importancia}     & Alta \\ \bottomrule
	\end{tabular}
	\caption{Caso de uso 18 - Visualizar productos}
\end{table}

\begin{table}[h]
	\centering
	\label{tabla:cu19}
	\begin{tabular}{@{}
		>{\columncolor[HTML]{FFFFFF}}p {.25\textwidth} p {.75\textwidth}@{}}
		\toprule
		\textbf{Caso de uso 19}   & \textbf{Filtrar productos} \\ \midrule
		\textbf{Versión}     & 2.0 \\ \midrule
		\textbf{Requisitos}	&  RF-4.6 \\ \midrule
		\textbf{Descripción}     & Permite filtrar los productos en función de los parámetros: cas, nombre, fórmula, nºeinecs y nºec. \\ \midrule
		\textbf{Precondiciones}  & 
		\begin{compactitem}
			\item El usuario se ha logueado en la aplicación con rol de Gestor de inventario.
		\end{compactitem}
		 \\ \midrule
		\textbf{Acciones} & 
		El usuario rellena o selecciona un valor de los desplegables en la columna de filtrado de la izquierda. En cuanto inserte un carácter el filtrado se producirá automáticamente sin pulsar ningún botón. También puede borrar todos los filtros pulsando en el botón Reiniciar filtro.
		\\ \midrule
		\textbf{Postcondiciones} & La aplicación mostrará en la lista únicamente los productos que coincidan con los parámetros de la búsqueda introducidos. \\ \midrule
		\textbf{Excepciones} & - \\ \midrule
		\textbf{Importancia} & Baja \\ \bottomrule
	\end{tabular}
	\caption{Caso de uso 19 - Filtrar productos}
\end{table}

\begin{table}[h]
	\centering
	\label{tabla:cu20}
	\begin{tabular}{@{}
		>{\columncolor[HTML]{FFFFFF}}p {.25\textwidth} p {.75\textwidth}@{}}
		\toprule
		\textbf{Caso de uso 20}   & \textbf{Insertar productos} \\ \midrule
		\textbf{Versión}     & 2.0 \\ \midrule
		\textbf{Requisitos}	&  RF-4.7 \\ \midrule
		\textbf{Descripción}     & Permite añadir un nuevo producto. \\ \midrule
		\textbf{Precondiciones}  & 
		\begin{compactitem}
			\item El usuario se ha logueado en la aplicación con rol de Gestor de inventario. 
		\end{compactitem}
		 \\ \midrule
		\textbf{Acciones} & 
		El usuario tiene que clickar sobre el botón de Nuevo producto. Después rellenará el cas, el nombre, la fórmula, pureza, peso molecular, f. desarrollada, nº einecs, nºec, precauciones, msds, peligros, prudencias y pictogramas y pulsará en el botón confirmar del modal.
		\\ \midrule
		\textbf{Postcondiciones} & La aplicación volverá a mostrar la lista con el producto añadido. \\ \midrule
		\textbf{Excepciones} & Si el producto ya existe, algún dato es inválido o no ha introducido uno de los campos obligatorios el sistema mostrará un mensaje de error. \\ \midrule
		\textbf{Importancia}     & Alta \\ \bottomrule
	\end{tabular}
	\caption{Caso de uso 20 - Insertar productos}
\end{table}

\begin{table}[h]
	\centering
	\label{tabla:cu21}
	\begin{tabular}{@{}
		>{\columncolor[HTML]{FFFFFF}}p {.25\textwidth} p {.75\textwidth}@{}}
		\toprule
		\textbf{Caso de uso 21}   & \textbf{Editar productos} \\ \midrule
		\textbf{Versión}     & 2.0 \\ \midrule
		\textbf{Requisitos}	&  RF-2.4, RF-2.6 \\ \midrule
		\textbf{Descripción}     & Permite al administrador editar los productos existentes. \\ \midrule
		\textbf{Precondiciones}  & 
		\begin{compactitem}
			\item El usuario se ha logueado en la aplicación con rol de Gestor de inventario. 
		\end{compactitem}
		 \\ \midrule
		\textbf{Acciones} & 
		El usuario tiene que clickar sobre el icono de la derecha de la fila del usuario que quiera editar. Después rellenará el cas, el nombre, la fórmula, pureza, peso molecular, f. desarrollada, nº einecs, nºec, precauciones, msds, peligros, prudencias y pictogramas y pulsará en el botón de confirmar del modal. 
		\\ \midrule
		\textbf{Postcondiciones} & La aplicación volverá a mostrar la lista con el producto editado. \\ \midrule
		\textbf{Excepciones} & Si el producto ya existe, algún dato es inválido o no ha introducido uno de los campos obligatorios el sistema mostrará un mensaje de error. \\ \midrule
		\textbf{Importancia}     & Media \\ \bottomrule
	\end{tabular}
	\caption{Caso de uso 21 - Editar productos}
\end{table}

%Gestion de ubicaciones -----------------------------
\begin{table}[h]
	\centering
	\label{tabla:cu22}
	\begin{tabular}{@{}
		>{\columncolor[HTML]{FFFFFF}}p {.25\textwidth} p {.75\textwidth}@{}}
		\toprule
		\textbf{Caso de uso 22}   & \textbf{Visualizar ubicaciones} \\ \midrule
		\textbf{Versión}     & 2.0 \\ \midrule
		\textbf{Requisitos}	&  RF-4.1, RF-4.5\\ \midrule
		\textbf{Descripción}     & Permite al administrador visualizar las ubicaciones existentes. \\ \midrule
		\textbf{Precondiciones}  & 
		\begin{compactitem}
			\item El usuario se ha logueado en la aplicación con rol de Gestor de inventario. 
			\item La ubicación ha de estar en la misma área a la que pertenece el usuario
		\end{compactitem}
		 \\ \midrule
		\textbf{Acciones} & 
		\begin{compactitem}
			\item El gestor se dirige a la página de ubicaciones 
		\end{compactitem}
		\\ \midrule
		\textbf{Postcondiciones} & -  \\ \midrule
		\textbf{Excepciones} &   - \\ \midrule
		\textbf{Importancia}     & Alta \\ \bottomrule
	\end{tabular}
	\caption{Caso de uso 22 - Visualizar ubicaciones}
\end{table}

\begin{table}[h]
	\centering
	\label{tabla:cu23}
	\begin{tabular}{@{}
		>{\columncolor[HTML]{FFFFFF}}p {.25\textwidth} p {.75\textwidth}@{}}
		\toprule
		\textbf{Caso de uso 23}   & \textbf{Filtrar ubicaciones} \\ \midrule
		\textbf{Versión}     & 2.0 \\ \midrule
		\textbf{Requisitos}	&  RF-4.6 \\ \midrule
		\textbf{Descripción}     & Permite filtrar las ubicaciones en función de los parámetros: nombre, departamento, área, centro y ubicación oculta. \\ \midrule
		\textbf{Precondiciones}  & 
		\begin{compactitem}
			\item El usuario se ha logueado en la aplicación con rol de Gestor de inventario.
		\end{compactitem}
		 \\ \midrule
		\textbf{Acciones} & 
		El usuario rellena o selecciona un valor de los desplegables en la columna de filtrado de la izquierda. En cuanto inserte un carácter el filtrado se producirá automáticamente sin pulsar ningún botón. También puede borrar todos los filtros pulsando en el botón Reiniciar filtro.
		\\ \midrule
		\textbf{Postcondiciones} & La aplicación mostrará en la lista únicamente las ubicaciones que coincidan con los parámetros de la búsqueda introducidos. \\ \midrule
		\textbf{Excepciones} & - \\ \midrule
		\textbf{Importancia} & Baja \\ \bottomrule
	\end{tabular}
	\caption{Caso de uso 23 - Filtrar ubicaciones}
\end{table}

\begin{table}[h]
	\centering
	\label{tabla:cu24}
	\begin{tabular}{@{}
		>{\columncolor[HTML]{FFFFFF}}p {.25\textwidth} p {.75\textwidth}@{}}
		\toprule
		\textbf{Caso de uso 24}   & \textbf{Insertar ubicaciones} \\ \midrule
		\textbf{Versión}     & 2.0 \\ \midrule
		\textbf{Requisitos}	&  RF-4.7 \\ \midrule
		\textbf{Descripción}     & Permite añadir una nueva ubicación. \\ \midrule
		\textbf{Precondiciones}  & 
		\begin{compactitem}
			\item El usuario se ha logueado en la aplicación con rol de Gestor de inventario. 
		\end{compactitem}
		 \\ \midrule
		\textbf{Acciones} & 
		El usuario tiene que clickar sobre el botón de Nueva ubicación. Después rellenará el nombre, el área el centro, si es oculta o no y pulsará en el botón confirmar del modal.
		\\ \midrule
		\textbf{Postcondiciones} & La aplicación volverá a mostrar la lista con la ubicación añadida. \\ \midrule
		\textbf{Excepciones} & Si la ubicación ya existe, algún dato es inválido o no ha introducido uno de los campos obligatorios el sistema mostrará un mensaje de error. \\ \midrule
		\textbf{Importancia}     & Alta \\ \bottomrule
	\end{tabular}
	\caption{Caso de uso 24 - Insertar ubicaciones}
\end{table}

\begin{table}[h]
	\centering
	\label{tabla:cu25}
	\begin{tabular}{@{}
		>{\columncolor[HTML]{FFFFFF}}p {.25\textwidth} p {.75\textwidth}@{}}
		\toprule
		\textbf{Caso de uso 25}   & \textbf{Editar ubicaciones} \\ \midrule
		\textbf{Versión}     & 2.0 \\ \midrule
		\textbf{Requisitos}	&  RF-2.4, RF-2.6 \\ \midrule
		\textbf{Descripción}     & Permite al administrador editar las ubicaciones existentes. \\ \midrule
		\textbf{Precondiciones}  & 
		\begin{compactitem}
			\item El usuario se ha logueado en la aplicación con rol de Gestor de inventario. 
		\end{compactitem}
		 \\ \midrule
		\textbf{Acciones} & 
		El usuario tiene que clickar sobre el icono de la derecha de la fila del usuario que quiera editar. Después rellenará el nombre, el área el centro, si es oculta o no y pulsará en el botón de confirmar del modal. 
		\\ \midrule
		\textbf{Postcondiciones} & La aplicación volverá a mostrar la lista con la ubicación editada. \\ \midrule
		\textbf{Excepciones} & Si la ubicación ya existe, algún dato es inválido o no ha introducido uno de los campos obligatorios el sistema mostrará un mensaje de error. \\ \midrule
		\textbf{Importancia}     & Media \\ \bottomrule
	\end{tabular}
	\caption{Caso de uso 25 - Editar ubicaciones}
\end{table}

%Gestion de calidades -----------------------------
\begin{table}[h]
	\centering
	\label{tabla:cu26}
	\begin{tabular}{@{}
		>{\columncolor[HTML]{FFFFFF}}p {.25\textwidth} p {.75\textwidth}@{}}
		\toprule
		\textbf{Caso de uso 26}   & \textbf{Visualizar calidades} \\ \midrule
		\textbf{Versión}     & 2.0 \\ \midrule
		\textbf{Requisitos}	&  RF-4.1, RF-4.5\\ \midrule
		\textbf{Descripción}     & Permite al administrador visualizar las calidades existentes. \\ \midrule
		\textbf{Precondiciones}  & 
		\begin{compactitem}
			\item El usuario se ha logueado en la aplicación con rol de Gestor de inventario. 
		\end{compactitem}
		 \\ \midrule
		\textbf{Acciones} & 
		\begin{compactitem}
			\item El gestor se dirige a la página de calidades 
		\end{compactitem}
		\\ \midrule
		\textbf{Postcondiciones} & -  \\ \midrule
		\textbf{Excepciones} &   - \\ \midrule
		\textbf{Importancia}     & Alta \\ \bottomrule
	\end{tabular}
	\caption{Caso de uso 26 - Visualizar calidades}
\end{table}

\begin{table}[h]
	\centering
	\label{tabla:cu27}
	\begin{tabular}{@{}
		>{\columncolor[HTML]{FFFFFF}}p {.25\textwidth} p {.75\textwidth}@{}}
		\toprule
		\textbf{Caso de uso 27}   & \textbf{Filtrar calidades} \\ \midrule
		\textbf{Versión}     & 2.0 \\ \midrule
		\textbf{Requisitos}	&  RF-4.6 \\ \midrule
		\textbf{Descripción}     & Permite filtrar las calidades en función del nombre. \\ \midrule
		\textbf{Precondiciones}  & 
		\begin{compactitem}
			\item El usuario se ha logueado en la aplicación con rol de Gestor de inventario.
		\end{compactitem}
		 \\ \midrule
		\textbf{Acciones} & 
		El usuario rellena el nombre en la columna de filtrado de la izquierda. En cuanto inserte un carácter el filtrado se producirá automáticamente sin pulsar ningún botón. También puede borrar todos los filtros pulsando en el botón Reiniciar filtro.
		\\ \midrule
		\textbf{Postcondiciones} & La aplicación mostrará en la lista únicamente las calidades que coincidan con los parámetros de la búsqueda introducidos. \\ \midrule
		\textbf{Excepciones} & - \\ \midrule
		\textbf{Importancia} & Baja \\ \bottomrule
	\end{tabular}
	\caption{Caso de uso 27 - Filtrar calidades}
\end{table}

\begin{table}[h]
	\centering
	\label{tabla:cu28}
	\begin{tabular}{@{}
		>{\columncolor[HTML]{FFFFFF}}p {.25\textwidth} p {.75\textwidth}@{}}
		\toprule
		\textbf{Caso de uso 28}   & \textbf{Insertar calidades} \\ \midrule
		\textbf{Versión}     & 2.0 \\ \midrule
		\textbf{Requisitos}	&  RF-4.7 \\ \midrule
		\textbf{Descripción}     & Permite añadir una nueva calidad. \\ \midrule
		\textbf{Precondiciones}  & 
		\begin{compactitem}
			\item El usuario se ha logueado en la aplicación con rol de Gestor de inventario. 
		\end{compactitem}
		 \\ \midrule
		\textbf{Acciones} & 
		El usuario tiene que clickar sobre el botón de Nueva calidad. Después rellenará el nombre y pulsará en el botón confirmar del modal.
		\\ \midrule
		\textbf{Postcondiciones} & La aplicación volverá a mostrar la lista con la calidad añadida. \\ \midrule
		\textbf{Excepciones} & Si la calidad ya existe, algún dato es inválido o no ha introducido uno de los campos obligatorios el sistema mostrará un mensaje de error. \\ \midrule
		\textbf{Importancia}     & Alta \\ \bottomrule
	\end{tabular}
	\caption{Caso de uso 28 - Insertar calidades}
\end{table}

\begin{table}[h]
	\centering
	\label{tabla:cu29}
	\begin{tabular}{@{}
		>{\columncolor[HTML]{FFFFFF}}p {.25\textwidth} p {.75\textwidth}@{}}
		\toprule
		\textbf{Caso de uso 29}   & \textbf{Editar calidades} \\ \midrule
		\textbf{Versión}     & 2.0 \\ \midrule
		\textbf{Requisitos}	&  RF-2.4, RF-2.6 \\ \midrule
		\textbf{Descripción}     & Permite al administrador editar las calidades existentes. \\ \midrule
		\textbf{Precondiciones}  & 
		\begin{compactitem}
			\item El usuario se ha logueado en la aplicación con rol de Gestor de inventario. 
		\end{compactitem}
		 \\ \midrule
		\textbf{Acciones} & 
		El usuario tiene que clickar sobre el icono de la derecha de la fila del usuario que quiera editar. Después rellenará el nombre y pulsará en el botón de confirmar del modal. 
		\\ \midrule
		\textbf{Postcondiciones} & La aplicación volverá a mostrar la lista con la calidad editada. \\ \midrule
		\textbf{Excepciones} & Si la calidad ya existe, algún dato es inválido o no ha introducido uno de los campos obligatorios el sistema mostrará un mensaje de error. \\ \midrule
		\textbf{Importancia}     & Media \\ \bottomrule
	\end{tabular}
	\caption{Caso de uso 29 - Editar calidades}
\end{table}

%Gestion de marcas -----------------------------
\begin{table}[h]
	\centering
	\label{tabla:cu30}
	\begin{tabular}{@{}
		>{\columncolor[HTML]{FFFFFF}}p {.25\textwidth} p {.75\textwidth}@{}}
		\toprule
		\textbf{Caso de uso 30}   & \textbf{Visualizar marcas} \\ \midrule
		\textbf{Versión}     & 2.0 \\ \midrule
		\textbf{Requisitos}	&  RF-4.1, RF-4.5\\ \midrule
		\textbf{Descripción}     & Permite al administrador visualizar las marcas existentes. \\ \midrule
		\textbf{Precondiciones}  & 
		\begin{compactitem}
			\item El usuario se ha logueado en la aplicación con rol de Gestor de inventario. 
		\end{compactitem}
		 \\ \midrule
		\textbf{Acciones} & 
		\begin{compactitem}
			\item El gestor se dirige a la página de marcas 
		\end{compactitem}
		\\ \midrule
		\textbf{Postcondiciones} & -  \\ \midrule
		\textbf{Excepciones} &   - \\ \midrule
		\textbf{Importancia}     & Alta \\ \bottomrule
	\end{tabular}
	\caption{Caso de uso 30 - Visualizar marcas}
\end{table}

\begin{table}[h]
	\centering
	\label{tabla:cu31}
	\begin{tabular}{@{}
		>{\columncolor[HTML]{FFFFFF}}p {.25\textwidth} p {.75\textwidth}@{}}
		\toprule
		\textbf{Caso de uso 31}   & \textbf{Filtrar marcas} \\ \midrule
		\textbf{Versión}     & 2.0 \\ \midrule
		\textbf{Requisitos}	&  RF-4.6 \\ \midrule
		\textbf{Descripción}     & Permite filtrar las marcas en función del nombre. \\ \midrule
		\textbf{Precondiciones}  & 
		\begin{compactitem}
			\item El usuario se ha logueado en la aplicación con rol de Gestor de inventario.
		\end{compactitem}
		 \\ \midrule
		\textbf{Acciones} & 
		El usuario rellena el nombre en la columna de filtrado de la izquierda. En cuanto inserte un carácter el filtrado se producirá automáticamente sin pulsar ningún botón. También puede borrar todos los filtros pulsando en el botón Reiniciar filtro.
		\\ \midrule
		\textbf{Postcondiciones} & La aplicación mostrará en la lista únicamente las marcas que coincidan con los parámetros de la búsqueda introducidos. \\ \midrule
		\textbf{Excepciones} & - \\ \midrule
		\textbf{Importancia} & Baja \\ \bottomrule
	\end{tabular}
	\caption{Caso de uso 31 - Filtrar marcas}
\end{table}

\begin{table}[h]
	\centering
	\label{tabla:cu32}
	\begin{tabular}{@{}
		>{\columncolor[HTML]{FFFFFF}}p {.25\textwidth} p {.75\textwidth}@{}}
		\toprule
		\textbf{Caso de uso 32}   & \textbf{Insertar marcas} \\ \midrule
		\textbf{Versión}     & 2.0 \\ \midrule
		\textbf{Requisitos}	&  RF-4.7 \\ \midrule
		\textbf{Descripción}     & Permite añadir una nueva marca. \\ \midrule
		\textbf{Precondiciones}  & 
		\begin{compactitem}
			\item El usuario se ha logueado en la aplicación con rol de Gestor de inventario. 
		\end{compactitem}
		 \\ \midrule
		\textbf{Acciones} & 
		El usuario tiene que clickar sobre el botón de Nueva marca. Después rellenará el nombre, teléfono, dirección y proveedores y pulsará en el botón confirmar del modal.
		\\ \midrule
		\textbf{Postcondiciones} & La aplicación volverá a mostrar la lista con la marca añadida. \\ \midrule
		\textbf{Excepciones} & Si la marca ya existe, algún dato es inválido o no ha introducido uno de los campos obligatorios el sistema mostrará un mensaje de error. \\ \midrule
		\textbf{Importancia}     & Alta \\ \bottomrule
	\end{tabular}
	\caption{Caso de uso 32 - Insertar marcas}
\end{table}

\begin{table}[h]
	\centering
	\label{tabla:cu33}
	\begin{tabular}{@{}
		>{\columncolor[HTML]{FFFFFF}}p {.25\textwidth} p {.75\textwidth}@{}}
		\toprule
		\textbf{Caso de uso 33}   & \textbf{Editar marcas} \\ \midrule
		\textbf{Versión}     & 2.0 \\ \midrule
		\textbf{Requisitos}	&  RF-2.4, RF-2.6 \\ \midrule
		\textbf{Descripción}     & Permite al administrador editar las marcas existentes. \\ \midrule
		\textbf{Precondiciones}  & 
		\begin{compactitem}
			\item El usuario se ha logueado en la aplicación con rol de Gestor de inventario. 
		\end{compactitem}
		 \\ \midrule
		\textbf{Acciones} & 
		El usuario tiene que clickar sobre el icono de la derecha de la fila del usuario que quiera editar. Después rellenará el nombre, teléfono, dirección y proveedores y pulsará en el botón de confirmar del modal. 
		\\ \midrule
		\textbf{Postcondiciones} & La aplicación volverá a mostrar la lista con la marca editada. \\ \midrule
		\textbf{Excepciones} & Si la marca ya existe, algún dato es inválido o no ha introducido uno de los campos obligatorios el sistema mostrará un mensaje de error. \\ \midrule
		\textbf{Importancia}     & Media \\ \bottomrule
	\end{tabular}
	\caption{Caso de uso 33 - Editar marcas}
\end{table}

%Gestion de proveedores -----------------------------
\begin{table}[h]
	\centering
	\label{tabla:cu34}
	\begin{tabular}{@{}
		>{\columncolor[HTML]{FFFFFF}}p {.25\textwidth} p {.75\textwidth}@{}}
		\toprule
		\textbf{Caso de uso 34}   & \textbf{Visualizar proveedores} \\ \midrule
		\textbf{Versión}     & 2.0 \\ \midrule
		\textbf{Requisitos}	&  RF-4.1, RF-4.5\\ \midrule
		\textbf{Descripción}     & Permite al administrador visualizar los proveedores existentes. \\ \midrule
		\textbf{Precondiciones}  & 
		\begin{compactitem}
			\item El usuario se ha logueado en la aplicación con rol de Gestor de inventario. 
		\end{compactitem}
		 \\ \midrule
		\textbf{Acciones} & 
		\begin{compactitem}
			\item El gestor se dirige a la página de proveedores 
		\end{compactitem}
		\\ \midrule
		\textbf{Postcondiciones} & -  \\ \midrule
		\textbf{Excepciones} &   - \\ \midrule
		\textbf{Importancia}     & Alta \\ \bottomrule
	\end{tabular}
	\caption{Caso de uso 34 - Visualizar proveedores}
\end{table}

\begin{table}[h]
	\centering
	\label{tabla:cu35}
	\begin{tabular}{@{}
		>{\columncolor[HTML]{FFFFFF}}p {.25\textwidth} p {.75\textwidth}@{}}
		\toprule
		\textbf{Caso de uso 35}   & \textbf{Filtrar proveedores} \\ \midrule
		\textbf{Versión}     & 2.0 \\ \midrule
		\textbf{Requisitos}	&  RF-4.6 \\ \midrule
		\textbf{Descripción}     & Permite filtrar los proveedores en función del nombre. \\ \midrule
		\textbf{Precondiciones}  & 
		\begin{compactitem}
			\item El usuario se ha logueado en la aplicación con rol de Gestor de inventario.
		\end{compactitem}
		 \\ \midrule
		\textbf{Acciones} & 
		El usuario rellena el nombre en la columna de filtrado de la izquierda. En cuanto inserte un carácter el filtrado se producirá automáticamente sin pulsar ningún botón. También puede borrar todos los filtros pulsando en el botón Reiniciar filtro.
		\\ \midrule
		\textbf{Postcondiciones} & La aplicación mostrará en la lista únicamente los proveedores que coincidan con los parámetros de la búsqueda introducidos. \\ \midrule
		\textbf{Excepciones} & - \\ \midrule
		\textbf{Importancia} & Baja \\ \bottomrule
	\end{tabular}
	\caption{Caso de uso 35 - Filtrar proveedores}
\end{table}

\begin{table}[h]
	\centering
	\label{tabla:cu36}
	\begin{tabular}{@{}
		>{\columncolor[HTML]{FFFFFF}}p {.25\textwidth} p {.75\textwidth}@{}}
		\toprule
		\textbf{Caso de uso 36}   & \textbf{Insertar proveedores} \\ \midrule
		\textbf{Versión}     & 2.0 \\ \midrule
		\textbf{Requisitos}	&  RF-4.7 \\ \midrule
		\textbf{Descripción}     & Permite añadir una nueva proveedor. \\ \midrule
		\textbf{Precondiciones}  & 
		\begin{compactitem}
			\item El usuario se ha logueado en la aplicación con rol de Gestor de inventario. 
		\end{compactitem}
		 \\ \midrule
		\textbf{Acciones} & 
		El usuario tiene que clickar sobre el botón de Nueva proveedor. Después rellenará el nombre, teléfono, dirección y proveedores y pulsará en el botón confirmar del modal.
		\\ \midrule
		\textbf{Postcondiciones} & La aplicación volverá a mostrar la lista con el proveedor añadido. \\ \midrule
		\textbf{Excepciones} & Si el proveedor ya existe, algún dato es inválido o no ha introducido uno de los campos obligatorios el sistema mostrará un mensaje de error. \\ \midrule
		\textbf{Importancia}     & Alta \\ \bottomrule
	\end{tabular}
	\caption{Caso de uso 36 - Insertar proveedores}
\end{table}

\begin{table}[h]
	\centering
	\label{tabla:cu37}
	\begin{tabular}{@{}
		>{\columncolor[HTML]{FFFFFF}}p {.25\textwidth} p {.75\textwidth}@{}}
		\toprule
		\textbf{Caso de uso 37}   & \textbf{Editar proveedores} \\ \midrule
		\textbf{Versión}     & 2.0 \\ \midrule
		\textbf{Requisitos}	&  RF-2.4, RF-2.6 \\ \midrule
		\textbf{Descripción}     & Permite al administrador editar los proveedores existentes. \\ \midrule
		\textbf{Precondiciones}  & 
		\begin{compactitem}
			\item El usuario se ha logueado en la aplicación con rol de Gestor de inventario. 
		\end{compactitem}
		 \\ \midrule
		\textbf{Acciones} & 
		El usuario tiene que clickar sobre el icono de la derecha de la fila del usuario que quiera editar. Después rellenará el nombre, teléfono, dirección y proveedores y pulsará en el botón de confirmar del modal. 
		\\ \midrule
		\textbf{Postcondiciones} & La aplicación volverá a mostrar la lista con el proveedor editado. \\ \midrule
		\textbf{Excepciones} & Si el proveedor ya existe, algún dato es inválido o no ha introducido uno de los campos obligatorios el sistema mostrará un mensaje de error. \\ \midrule
		\textbf{Importancia}     & Media \\ \bottomrule
	\end{tabular}
	\caption{Caso de uso 37 - Editar proveedores}
\end{table}
\apendice{Especificación de diseño}

\section{Introducción}

\section{Diseño de datos}

\section{Diseño procedimental}

\section{Diseño arquitectónico}



\apendice{Documentación técnica de programación}

\section{Introducción}

En este apartado se explica todo lo que debe conocer el desarrollador a la hora de instalar el entorno de trabajo y poder continuar con el desarrollo de la aplicación.

\section{Estructura de directorios}

\dirtree{%
 .1 /.
 .2 design \desc{Aquí se encuentran los recursos destinados al diseño gráfico de la aplicación}.
 .2 doc \desc{Directorio de los documentos relativos a la documentación}.
 .2 MySQL \desc{Diagrama entidad-relacion y scripts para inicializar la base de datos}.
 .2 src \desc{Directorio donde encontramos la lógica de la aplicación}.
 .3 lib \desc{Directorio de librerías}.
 .3 main.
 .4 java \desc{Contenedor de los paquetes java}.
 .5 geslab \desc{Paquete padre}.
 .6 database \desc{Contiene las clases y objetos que conectan de alguna forma con la base de datos}.
 .7 admin \desc{Contiene las clases y objetos que se usan en la página de administración}.
 .7 user \desc{Contiene las clases y objetos que se usan en la página de mantenimiento}.
 .6 servlets \desc{Contiene todos los servlets}.
 .4 resources.
 .4 webapp \desc{Directorio que contiene toda la lógica de la parte web}.
 .5 css \desc{Hojas de estilo}.
 .5 images \desc{Imágenes utilizadas en la web}.
 .5 js \desc{Javascripts de cada archivo jsp}.
 .5 WEB-INF \desc{Contiene las vistas jsp}.
 .2 target \desc{Carpeta contenedora de los documentos de las dependencias del proyecto}.
}

\section{Manual del programador}

En este apartado se se explicará como configurar todo el entorno de trabajo para poder en un futuro si es necesario, continuar con el desarrollo de la aplicación.

El desarrollo del proyecto se realizó sobre un sistema operativo Windows, por lo que todos los programas a los que se hará referencia serán en su versión para este sistema operativo.

\subsection{MySQL y MySQL Workbench}

El primer paso para instalar MYSQL es descargar el software de su página oficial \href{https://dev.mysql.com/downloads/mysql/}{\textit{MySql}}\footnote{\textit{MySql}: \url{https://dev.mysql.com/downloads/mysql/}}. Desde la página buscaremos el MySQL Comunnity Server y después seleccionar la plataforma en la que lo vamos a instalar, en nuestro caso windows.

Una vez tengamos el ejecutable solo hay que abrirlo y seleccionar el modo de instalación. En nuestro caso selecciono Developer Default ya que quiero instalar tanto MySQL Server como MySql Workbech, un software que proporciona diversas opciones de administraciónde de bases de datos relacionales así como herramientas que nos permiten utilizar consultas SQL mediante un entorno visual.

En cierto punto de la instalación puede que el sistema nos diga que necesitamos el paquete de Microsoft Visual C++, solo tenemos que continuar para que se instale automáticamente.

Finalizada la instalación, solo nos faltaría establecer una configuración inicial antes de ejecutar los correspondientes servicios. En la primera ventana seleccionaremos \textit{Standalone MySQL Server/Classic MySQL Replication}. En Type and Networking dejamos todos los parámetros por defecto, a no ser que quisieramos establecer la conexión en otro puerto distinto al 3306. En caso de cambiar el puerto es importante recordarlo para futuras configuraciones.

Después debemos elegir una contraseña para poder conectarnos al servidor SQL. No es necesario que definamos un usuario para administrar la base de datos, ya que podemos utilizar el usuario por defecto \textit{root}.

Por último definimos un nombre para el servicio MySQL y las preferencias generales y pulsamos en execute para ejecutar las acciones y activar todos los servicios

\subsection{Java}

Para la instalación de Java en el equipo solamente tenemos que ir a la página de descarga de la propia web de \href{https://www.oracle.com/java/technologies/javase/jdk14-archive-downloads.html}{\textit{Oracle}}\footnote{\textit{JDK}: \url{https://www.oracle.com/java/technologies/javase/jdk14-archive-downloads.html}}. Desde ahí nos descargaremos el archivo \textit{Windows x64 Installer}. Ejecutamos el archivo y pulsamos en siguiente hasta completar la instalación. Para asegurarnos que se ha instalado podemos irnos a la ventana de comando y ejecutar \textit{java -version}

\subsection{Eclipse IDE}

Para la instalación de eclipse, nos iremos en primer lugar a su página para descargar el instalador de \href{https://www.eclipse.org/downloads/}{\textit{Eclipse IDE}}\footnote{\textit{Eclipse}: \url{https://www.eclipse.org/downloads/}}.

Una vez con el ejecutable, lo abrimos y seleccionamos \textit{Eclipse IDE for Enterprise Java Developers}. Dejamos las rutas por defecto y pulsamos en INSTALL. 

\section{Compilación, instalación y ejecución del proyecto}

Con todo ya instalado, para poder ejecutar nuestra aplicación deberemos de seguir los pasos que se indican a continuación.

\subsection{Inicializar base de datos}

Abrimos MySQL Workbench y nos conectamos a la instancia que anteriormente creamos. Una vez conectados, debemos de crear un nuevo schema con el nombre de \textbf{dbo}. Para ello en la barra de herramientas superior buscamos el icono pertinente y lo pulsamos.

Una vez creado el schema, abriremos el archivo sql que se encuentra en la carpeta \textit{MySQL} llamado \textit{script-create-tables} y lo ejecutaremos. 

Con esto, ya se nos habrán creado todas las tablas necesarias, junto con el usuario administrador y otros datos inicializados.

\subsection{Importar proyecto en Eclipse}

Para importar el proyecto en Eclipse, tan solo tenemos que abrir el programa y dirigirnos a \textit{File/Import}. Desde aquí buscaremos la opción de \textit{Existing Maven Projects} dentro del apartado de Maven.

desde ahí seleccionaremos la ruta de origen de nuestro proyecto, y seleccionaremos el proyecto una vez que aparezca en la lista de disponibles. Con esto ya tendremos nuestro proyecto importado junto con todas sus dependencias.

\subsection{Configurar conexión con la base de datos}

Por último, para que la conexión con la base de datos sea satisfactoria, tenemos que dirigirnos a la clase \textit{Conexion}, dentro del paquete \textit{geslab.database}. 

En los atributos de la clase tendremos que editar una serie de valores en función de los datos con los que configuramos la base de datos.

En la imagen \ref{fig:configurarConexion} se pueden ver los campos que habría que editar marcados en rojo en cursiva y entre asteriscos.

\imagen{configurarConexion}{Configurar conexión con la base de datos} 

\section{Pruebas del sistema}

Se han realizado dos tipos de pruebas. Por un lado se han llevado a cabo las\textbf{ pruebas de cada caso de uso} de una manera manual. Para mostrar sus resultados se ha utilizado una plantilla provista por la junta de Andalucía \cite{doc:pruebas}. En las siguientes tablas se pueden encontrar todos los casos de prueba para cada caso de uso.

%Gestion de centros -----------------------------
\begin{table}[h]
	\centering
	\label{tabla:prueba1}
	\begin{tabular}{@{}
		>{\columncolor[HTML]{FFFFFF}}p {.25\textwidth} p {.75\textwidth}@{}}
		\toprule
		\textbf{Caso de prueba 1} & \textbf{Visualizar centros} \\ \midrule
		\textbf{Descripción}     & Comprobar la correcta visualización de los centros \\ \midrule
		\textbf{Prerrequisitos}	&  Loguearse como administrador \\ \midrule
		\textbf{Pasos}  & 
		\begin{compactitem}
			\item  Desplazarse a la pestaña de centros
		\end{compactitem}
		 \\ \midrule
		\textbf{Resultado esperado} & Ver la lista de centros
		\\ \midrule
		\textbf{Resultado obtenido} & Visualizada lista de centros\\ \midrule
	\end{tabular}
	\caption{Caso de prueba 1 - Visualizar centros}
\end{table}

\begin{table}[h]
	\centering
	\label{tabla:prueba2}
	\begin{tabular}{@{}
		>{\columncolor[HTML]{FFFFFF}}p {.25\textwidth} p {.75\textwidth}@{}}
		\toprule
		\textbf{Caso de prueba 2}   & \textbf{Editar centros} \\ \midrule
		\textbf{Descripción}	& Comprobar la correcta actualización del centro \\ \midrule
		\textbf{Prerrequisitos} & Loguearse como administrador \\ \midrule
		\textbf{Pasos}  & 
		\begin{compactitem}
			\item Desplazarse a la pestaña de centros
			\item Pinchar el icono de la derecha del centro a editar
			\item Editar los campos 
		\end{compactitem}
		 \\ \midrule
		\textbf{Resultado esperado} & 
		Se muestra el centro con los nuevos datos introducidos
		\\ \midrule
		\textbf{Resultado obtenido} & El centro muestra los datos actualizados correctamente\\ \midrule
	\end{tabular}
	\caption{Caso de prueba 2 - Editar centros}
\end{table}

\begin{table}[h]
	\centering
	\label{tabla:prueba3}
	\begin{tabular}{@{}
		>{\columncolor[HTML]{FFFFFF}}p {.25\textwidth} p {.75\textwidth}@{}}
		\toprule
		\textbf{Caso de prueba 3}   & \textbf{Insertar centros} \\ \midrule
		\textbf{Descripción}	&  Comprobar la correcta inserción del centro \\ \midrule
		\textbf{Prerrequisitos}   & Loguearse como administrador \\ \midrule
		\textbf{Pasos}  & 
		\begin{compactitem}
			\item Desplazarse a la pestaña de centros
			\item Pinchar el botón de Nuevo centro
			\item Insertar los campos
			\item Pulsar en confirmar  
		\end{compactitem}
		 \\ \midrule
		\textbf{Resultado esperado} & 
		Se muestra el nuevo centro en la lista de centros
		\\ \midrule
		\textbf{Resultado obtenido} & El centro se muestra correctamente \\ \midrule
	\end{tabular}
	\caption{Caso de prueba 3 - Insertar centros}
\end{table}

%Gestion de departamentos -----------------------------
\begin{table}[h]
	\centering
	\label{tabla:prueba4}
	\begin{tabular}{@{}
		>{\columncolor[HTML]{FFFFFF}}p {.25\textwidth} p {.75\textwidth}@{}}
		\toprule
		\textbf{Caso de prueba 4}   & \textbf{Visualizar departamentos} \\ \midrule
		\textbf{Descripción}     & Comprobar la correcta visualización de los departamentos \\ \midrule
		\textbf{Prerrequisitos}	&  Loguearse como administrador \\ \midrule
		\textbf{Pasos}  & 
		\begin{compactitem}
			\item  Desplazarse a la pestaña de departamentos
		\end{compactitem}
		 \\ \midrule
		\textbf{Resultado esperado} & Ver la lista de departamentos
		\\ \midrule
		\textbf{Resultado obtenido} & Visualizada lista de departamentos\\ \midrule
	\end{tabular}
	\caption{Caso de prueba 4 - Visualizar departamentos}
\end{table}

\begin{table}[h]
	\centering
	\label{tabla:prueba5}
	\begin{tabular}{@{}
		>{\columncolor[HTML]{FFFFFF}}p {.25\textwidth} p {.75\textwidth}@{}}
		\toprule
		\textbf{Caso de prueba 5}   & \textbf{Editar departamentos} \\ \midrule
		\textbf{Descripción}	& Comprobar la correcta actualización del departamento \\ \midrule
		\textbf{Prerrequisitos} & Loguearse como administrador \\ \midrule
		\textbf{Pasos}  & 
		\begin{compactitem}
			\item Desplazarse a la pestaña de departamentos
			\item Pinchar el icono de la derecha del departamento a editar
			\item Editar los campos 
		\end{compactitem}
		 \\ \midrule
		\textbf{Resultado esperado} & 
		Se muestra el departamento con los nuevos datos introducidos
		\\ \midrule
		\textbf{Resultado obtenido} & El departamento muestra los datos actualizados correctamente\\ \midrule
	\end{tabular}
	\caption{Caso de prueba 5 - Editar departamentos}
\end{table}

\begin{table}[h]
	\centering
	\label{tabla:prueba6}
	\begin{tabular}{@{}
		>{\columncolor[HTML]{FFFFFF}}p {.25\textwidth} p {.75\textwidth}@{}}
		\toprule
		\textbf{Caso de prueba 6}   & \textbf{Insertar departamentos} \\ \midrule
		\textbf{Descripción}	&  Comprobar la correcta inserción del departamento \\ \midrule
		\textbf{Prerrequisitos}   & Loguearse como administrador \\ \midrule
		\textbf{Pasos}  & 
		\begin{compactitem}
			\item Desplazarse a la pestaña de departamentos
			\item Pinchar el botón de Nuevo departamento
			\item Insertar los campos
			\item Pulsar en confirmar  
		\end{compactitem}
		 \\ \midrule
		\textbf{Resultado esperado} & 
		Se muestra el nuevo departamento en la lista de departamentos
		\\ \midrule
		\textbf{Resultado obtenido} & El departamento se muestra correctamente \\ \midrule
	\end{tabular}
	\caption{Caso de prueba 6 - Insertar departamentos}
\end{table}


%Gestion de áreas -----------------------------
\begin{table}[h]
	\centering
	\label{tabla:prueba7}
	\begin{tabular}{@{}
		>{\columncolor[HTML]{FFFFFF}}p {.25\textwidth} p {.75\textwidth}@{}}
		\toprule
		\textbf{Caso de prueba 7}   & \textbf{Visualizar áreas} \\ \midrule
		\textbf{Descripción}     & Comprobar la correcta visualización de las áreas \\ \midrule
		\textbf{Prerrequisitos}	&  Loguearse como administrador \\ \midrule
		\textbf{Pasos}  & 
		\begin{compactitem}
			\item  Desplazarse a la pestaña de áreas
		\end{compactitem}
		 \\ \midrule
		\textbf{Resultado esperado} & Ver la lista de áreas
		\\ \midrule
		\textbf{Resultado obtenido} & Visualizada lista de áreas\\ \midrule
	\end{tabular}
	\caption{Caso de prueba 7 - Visualizar áreas}
\end{table}

\begin{table}[h]
	\centering
	\label{tabla:prueba8}
	\begin{tabular}{@{}
		>{\columncolor[HTML]{FFFFFF}}p {.25\textwidth} p {.75\textwidth}@{}}
		\toprule
		\textbf{Caso de prueba 8}   & \textbf{Editar áreas} \\ \midrule
		\textbf{Descripción}	& Comprobar la correcta actualización del área \\ \midrule
		\textbf{Prerrequisitos} & Loguearse como administrador \\ \midrule
		\textbf{Pasos}  & 
		\begin{compactitem}
			\item Desplazarse a la pestaña de áreas
			\item Pinchar el icono de la derecha del área a editar
			\item Editar los campos 
		\end{compactitem}
		 \\ \midrule
		\textbf{Resultado esperado} & 
		Se muestra el área con los nuevos datos introducidos
		\\ \midrule
		\textbf{Resultado obtenido} & El área muestra los datos actualizados correctamente\\ \midrule
	\end{tabular}
	\caption{Caso de prueba 8 - Editar áreas}
\end{table}

\begin{table}[h]
	\centering
	\label{tabla:prueba9}
	\begin{tabular}{@{}
		>{\columncolor[HTML]{FFFFFF}}p {.25\textwidth} p {.75\textwidth}@{}}
		\toprule
		\textbf{Caso de prueba 9}   & \textbf{Insertar áreas} \\ \midrule
		\textbf{Descripción}	&  Comprobar la correcta inserción del área \\ \midrule
		\textbf{Prerrequisitos}   & Loguearse como administrador \\ \midrule
		\textbf{Pasos}  & 
		\begin{compactitem}
			\item Desplazarse a la pestaña de áreas
			\item Pinchar el botón de Nuevo área
			\item Insertar los campos
			\item Pulsar en confirmar  
		\end{compactitem}
		 \\ \midrule
		\textbf{Resultado esperado} & 
		Se muestra el nuevo área en la lista de áreas
		\\ \midrule
		\textbf{Resultado obtenido} & El área se muestra correctamente \\ \midrule
	\end{tabular}
	\caption{Caso de prueba 9 - Insertar áreas}
\end{table}

%Gestion de usuarios -----------------------------
\begin{table}[h]
	\centering
	\label{tabla:prueba10}
	\begin{tabular}{@{}
		>{\columncolor[HTML]{FFFFFF}}p {.25\textwidth} p {.75\textwidth}@{}}
		\toprule
		\textbf{Caso de prueba 10}   & \textbf{Visualizar usuarios} \\ \midrule
		\textbf{Descripción}     & Comprobar la correcta visualización de los usuarios \\ \midrule
		\textbf{Prerrequisitos}	&  Loguearse como administrador \\ \midrule
		\textbf{Pasos}  & 
		\begin{compactitem}
			\item  Desplazarse a la pestaña de usuarios
		\end{compactitem}
		 \\ \midrule
		\textbf{Resultado esperado} & Ver la lista de usuarios
		\\ \midrule
		\textbf{Resultado obtenido} & Visualizada lista de usuarios\\ \midrule
	\end{tabular}
	\caption{Caso de prueba 10 - Visualizar usuarios}
\end{table}

\begin{table}[h]
	\centering
	\label{tabla:prueba11}
	\begin{tabular}{@{}
		>{\columncolor[HTML]{FFFFFF}}p {.25\textwidth} p {.75\textwidth}@{}}
		\toprule
		\textbf{Caso de prueba 11}   & \textbf{Editar usuarios} \\ \midrule
		\textbf{Descripción}	& Comprobar la correcta actualización del usuario \\ \midrule
		\textbf{Prerrequisitos} & Loguearse como administrador \\ \midrule
		\textbf{Pasos}  & 
		\begin{compactitem}
			\item Desplazarse a la pestaña de usuarios
			\item Pinchar el icono de la derecha del usuario a editar
			\item Editar los campos 
		\end{compactitem}
		 \\ \midrule
		\textbf{Resultado esperado} & 
		Se muestra el usuario con los nuevos datos introducidos
		\\ \midrule
		\textbf{Resultado obtenido} & El usuario muestra los datos actualizados correctamente\\ \midrule
	\end{tabular}
	\caption{Caso de prueba 11 - Editar usuarios}
\end{table}

\begin{table}[h]
	\centering
	\label{tabla:prueba12}
	\begin{tabular}{@{}
		>{\columncolor[HTML]{FFFFFF}}p {.25\textwidth} p {.75\textwidth}@{}}
		\toprule
		\textbf{Caso de prueba 12}   & \textbf{Insertar usuarios} \\ \midrule
		\textbf{Descripción}	&  Comprobar la correcta inserción del usuario \\ \midrule
		\textbf{Prerrequisitos}   & Loguearse como administrador \\ \midrule
		\textbf{Pasos}  & 
		\begin{compactitem}
			\item Desplazarse a la pestaña de usuarios
			\item Pinchar el botón de Nuevo usuario
			\item Insertar los campos
			\item Pulsar en confirmar  
		\end{compactitem}
		 \\ \midrule
		\textbf{Resultado esperado} & 
		Se muestra el nuevo usuario en la lista de usuarios
		\\ \midrule
		\textbf{Resultado obtenido} & El usuario se muestra correctamente \\ \midrule
	\end{tabular}
	\caption{Caso de prueba 12 - Insertar usuarios}
\end{table}


%Gestion de existencias -----------------------------
\begin{table}[h]
	\centering
	\label{tabla:prueba13}
	\begin{tabular}{@{}
		>{\columncolor[HTML]{FFFFFF}}p {.25\textwidth} p {.75\textwidth}@{}}
		\toprule
		\textbf{Caso de prueba 13}   & \textbf{Visualizar existencias} \\ \midrule
		\textbf{Descripción}     & Comprobar la correcta visualización de las fichas \\ \midrule
		\textbf{Prerrequisitos}	&  Loguearse como gestor de inventario \\ \midrule
		\textbf{Pasos}  & 
		\begin{compactitem}
			\item  Desplazarse a la pestaña de existencias
		\end{compactitem}
		 \\ \midrule
		\textbf{Resultado esperado} & Ver la lista de fichas
		\\ \midrule
		\textbf{Resultado obtenido} & Visualizada lista de existencias\\ \midrule
	\end{tabular}
	\caption{Caso de prueba 13 - Visualizar existencias}
\end{table}

\begin{table}[h]
	\centering
	\label{tabla:prueba14}
	\begin{tabular}{@{}
		>{\columncolor[HTML]{FFFFFF}}p {.25\textwidth} p {.75\textwidth}@{}}
		\toprule
		\textbf{Caso de prueba 14}   & \textbf{Filtrar existencias} \\ \midrule
		\textbf{Descripción}	&  Comprobar el sistema de filtrado de fichas \\ \midrule
		\textbf{Prerrequisitos} & Loguearse como gestor de inventario\\ \midrule
		\textbf{Pasos}  & 
		\begin{compactitem}
			\item Desplazarse a la ventana de existencias desde el menu
			\item Establecer parámetros en la columna de filtros
		\end{compactitem}
		 \\ \midrule
		\textbf{Resultado esperado} & 
		La lista de fichas solo muestra las filas que cumplen los parámetros indicados
		\\ \midrule
		\textbf{Resultado obtenido} & Solo se pueden ver fichas con los parámetros filtrados. \\ \midrule
	\end{tabular}
	\caption{Caso de prueba 14 - Filtrar existencias}
\end{table}

\begin{table}[h]
	\centering
	\label{tabla:prueba15}
	\begin{tabular}{@{}
		>{\columncolor[HTML]{FFFFFF}}p {.25\textwidth} p {.75\textwidth}@{}}
		\toprule
		\textbf{Caso de prueba 15}   & \textbf{Insertar fichas} \\ \midrule
		\textbf{Descripción}	&  Comprobar la correcta inserción del centro \\ \midrule
		\textbf{Prerrequisitos}   & Loguearse como gestor de inventario \\ \midrule
		\textbf{Pasos}  & 
		\begin{compactitem}
			\item Desplazarse a la pestaña de existencias
			\item Pinchar el botón de Nuevo centro
			\item Insertar los campos
			\item Pulsar en confirmar  
		\end{compactitem}
		 \\ \midrule
		\textbf{Resultado esperado} & 
		Se muestra el nuevo centro en la lista de existencias
		\\ \midrule
		\textbf{Resultado obtenido} & El centro se muestra correctamente \\ \midrule
	\end{tabular}
	\caption{Caso de prueba 15 - Insertar fichas}
\end{table}

\begin{table}[h]
	\centering
	\label{tabla:prueba16}
	\begin{tabular}{@{}
		>{\columncolor[HTML]{FFFFFF}}p {.25\textwidth} p {.75\textwidth}@{}}
		\toprule
		\textbf{Caso de prueba 16}   & \textbf{Añadir entradas} \\ \midrule
		\textbf{Descripción}	&  Comprobar el incremento del stock al añadir existencias \\ \midrule
		\textbf{Prerrequisitos} & Loguearse como gestor de inventario\\ \midrule
		\textbf{Pasos}  & 
		\begin{compactitem}
			\item Desplazarse a la ventana de existencias desde el menu
			\item Pinchar en el botón con un + a la derecha de la ficha
			\item Rellenar los datos solicitados
		\end{compactitem}
		 \\ \midrule
		\textbf{Resultado esperado} & 
		El stock a aumentado y la entrada se muestra en la lista de entradas
		\\ \midrule
		\textbf{Resultado obtenido} & El stock se ha incrementado correctamente, y la entrada aparece en la lista de entradas de la ficha \\ \midrule
	\end{tabular}
	\caption{Caso de prueba 16 - Añadir entradas}
\end{table}

\begin{table}[h]
	\centering
	\label{tabla:prueba17}
	\begin{tabular}{@{}
		>{\columncolor[HTML]{FFFFFF}}p {.25\textwidth} p {.75\textwidth}@{}}
		\toprule
		\textbf{Caso de prueba 17}   & \textbf{Añadir salidas} \\ \midrule
		\textbf{Descripción}	&  Comprobar el incremento del stock al añadir existencias \\ \midrule
		\textbf{Prerrequisitos} & Loguearse como gestor de inventario\\ \midrule
		\textbf{Pasos}  & 
		\begin{compactitem}
			\item Desplazarse a la ventana de existencias desde el menu
			\item Pinchar en el botón con un - a la derecha de la ficha
			\item Rellenar los datos solicitados
		\end{compactitem}
		 \\ \midrule
		\textbf{Resultado esperado} & 
		El stock a disminuido y la salida se muestra en la lista de salidas
		\\ \midrule
		\textbf{Resultado obtenido} & El stock se ha disminuido correctamente, y la salida aparece en la lista de salidas de la ficha \\ \midrule
	\end{tabular}
	\caption{Caso de prueba 17 - Añadir salidas}
\end{table}


%Gestion de productos -----------------------------
\begin{table}[h]
	\centering
	\label{tabla:prueba18}
	\begin{tabular}{@{}
		>{\columncolor[HTML]{FFFFFF}}p {.25\textwidth} p {.75\textwidth}@{}}
		\toprule
		\textbf{Caso de prueba 18}   & \textbf{Visualizar productos} \\ \midrule
		\textbf{Descripción}     & Comprobar la correcta visualización de los productos \\ \midrule
		\textbf{Prerrequisitos}	&  Loguearse como gestor de inventario \\ \midrule
		\textbf{Pasos}  & 
		\begin{compactitem}
			\item  Desplazarse a la pestaña de productos
		\end{compactitem}
		 \\ \midrule
		\textbf{Resultado esperado} & Ver la lista de productos
		\\ \midrule
		\textbf{Resultado obtenido} & Visualizada lista de productos\\ \midrule
	\end{tabular}
	\caption{Caso de prueba 18 - Visualizar productos}
\end{table}

\begin{table}[h]
	\centering
	\label{tabla:prueba19}
	\begin{tabular}{@{}
		>{\columncolor[HTML]{FFFFFF}}p {.25\textwidth} p {.75\textwidth}@{}}
		\toprule
		\textbf{Caso de prueba 19}   & \textbf{Filtrar productos} \\ \midrule
		\textbf{Descripción}	&  Comprobar el sistema de filtrado de productos \\ \midrule
		\textbf{Prerrequisitos} & Loguearse como gestor de inventario\\ \midrule
		\textbf{Pasos}  & 
		\begin{compactitem}
			\item Desplazarse a la ventana de productos desde el menu
			\item Establecer parámetros en la columna de filtros
		\end{compactitem}
		 \\ \midrule
		\textbf{Resultado esperado} & 
		La lista de productos solo muestra las filas que cumplen los parámetros indicados
		\\ \midrule
		\textbf{Resultado obtenido} & Solo se pueden ver productos con los parámetros filtrados. \\ \midrule
	\end{tabular}
	\caption{Caso de prueba 19 - Filtrar productos}
\end{table}

\begin{table}[h]
	\centering
	\label{tabla:prueba20}
	\begin{tabular}{@{}
		>{\columncolor[HTML]{FFFFFF}}p {.25\textwidth} p {.75\textwidth}@{}}
		\toprule
		\textbf{Caso de prueba 20}   & \textbf{Insertar productos} \\ \midrule
		\textbf{Descripción}	&  Comprobar la correcta inserción del producto \\ \midrule
		\textbf{Prerrequisitos}   & Loguearse como gestor de inventario \\ \midrule
		\textbf{Pasos}  & 
		\begin{compactitem}
			\item Desplazarse a la pestaña de productos
			\item Pinchar el botón de Nuevo producto
			\item Insertar los campos
			\item Pulsar en confirmar  
		\end{compactitem}
		 \\ \midrule
		\textbf{Resultado esperado} & 
		Se muestra el nuevo producto en la lista de productos
		\\ \midrule
		\textbf{Resultado obtenido} & El producto se muestra correctamente \\ \midrule
	\end{tabular}
	\caption{Caso de prueba 20 - Insertar productos}
\end{table}

\begin{table}[h]
	\centering
	\label{tabla:prueba21}
	\begin{tabular}{@{}
		>{\columncolor[HTML]{FFFFFF}}p {.25\textwidth} p {.75\textwidth}@{}}
		\toprule
		\textbf{Caso de prueba 21}   & \textbf{Editar productos} \\ \midrule
		\textbf{Descripción}	& Comprobar la correcta actualización del producto \\ \midrule
		\textbf{Prerrequisitos} & Loguearse como gestor de inventario \\ \midrule
		\textbf{Pasos}  & 
		\begin{compactitem}
			\item Desplazarse a la pestaña de productos
			\item Pinchar el icono de la derecha del producto a editar
			\item Editar los campos 
		\end{compactitem}
		 \\ \midrule
		\textbf{Resultado esperado} & 
		Se muestra el producto con los nuevos datos introducidos
		\\ \midrule
		\textbf{Resultado obtenido} & El producto muestra los datos actualizados correctamente\\ \midrule
	\end{tabular}
	\caption{Caso de prueba 21 - Editar productos}
\end{table}

%Gestion de ubicaciones -----------------------------
\begin{table}[h]
	\centering
	\label{tabla:prueba22}
	\begin{tabular}{@{}
		>{\columncolor[HTML]{FFFFFF}}p {.25\textwidth} p {.75\textwidth}@{}}
		\toprule
		\textbf{Caso de prueba 22}   & \textbf{Visualizar ubicaciones} \\ \midrule
		\textbf{Descripción}     & Comprobar la correcta visualización de las ubicaciones \\ \midrule
		\textbf{Prerrequisitos}	&  Loguearse como gestor de inventario \\ \midrule
		\textbf{Pasos}  & 
		\begin{compactitem}
			\item  Desplazarse a la pestaña de ubicaciones
		\end{compactitem}
		 \\ \midrule
		\textbf{Resultado esperado} & Ver la lista de ubicaciones
		\\ \midrule
		\textbf{Resultado obtenido} & Visualizada lista de ubicaciones\\ \midrule
	\end{tabular}
	\caption{Caso de prueba 22 - Visualizar ubicaciones}
\end{table}

\begin{table}[h]
	\centering
	\label{tabla:prueba23}
	\begin{tabular}{@{}
		>{\columncolor[HTML]{FFFFFF}}p {.25\textwidth} p {.75\textwidth}@{}}
		\toprule
		\textbf{Caso de prueba 23}   & \textbf{Filtrar ubicaciones} \\ \midrule
		\textbf{Descripción}	&  Comprobar el sistema de filtrado de ubicaciones \\ \midrule
		\textbf{Prerrequisitos} & Loguearse como gestor de inventario\\ \midrule
		\textbf{Pasos}  & 
		\begin{compactitem}
			\item Desplazarse a la ventana de ubicaciones desde el menu
			\item Establecer parámetros en la columna de filtros
		\end{compactitem}
		 \\ \midrule
		\textbf{Resultado esperado} & 
		La lista de ubicaciones solo muestra las filas que cumplen los parámetros indicados
		\\ \midrule
		\textbf{Resultado obtenido} & Solo se pueden ver ubicaciones con los parámetros filtrados. \\ \midrule
	\end{tabular}
	\caption{Caso de prueba 23 - Filtrar ubicaciones}
\end{table}

\begin{table}[h]
	\centering
	\label{tabla:prueba24}
	\begin{tabular}{@{}
		>{\columncolor[HTML]{FFFFFF}}p {.25\textwidth} p {.75\textwidth}@{}}
		\toprule
		\textbf{Caso de prueba 24}   & \textbf{Insertar ubicaciones} \\ \midrule
		\textbf{Descripción}	&  Comprobar la correcta inserción de la ubicación\\ \midrule
		\textbf{Prerrequisitos}   & Loguearse como gestor de inventario \\ \midrule
		\textbf{Pasos}  & 
		\begin{compactitem}
			\item Desplazarse a la pestaña de ubicaciones
			\item Pinchar el botón de Nueva ubicación
			\item Insertar los campos
			\item Pulsar en confirmar  
		\end{compactitem}
		 \\ \midrule
		\textbf{Resultado esperado} & 
		Se muestra la nueva ubicación en la lista de ubicaciones
		\\ \midrule
		\textbf{Resultado obtenido} & La ubicación se muestra correctamente \\ \midrule
	\end{tabular}
	\caption{Caso de prueba 24 - Insertar ubicaciones}
\end{table}

\begin{table}[h]
	\centering
	\label{tabla:prueba25}
	\begin{tabular}{@{}
		>{\columncolor[HTML]{FFFFFF}}p {.25\textwidth} p {.75\textwidth}@{}}
		\toprule
		\textbf{Caso de prueba 25}   & \textbf{Editar ubicaciones} \\ \midrule
		\textbf{Descripción}	& Comprobar la correcta actualización de la ubicación \\ \midrule
		\textbf{Prerrequisitos} & Loguearse como gestor de inventario \\ \midrule
		\textbf{Pasos}  & 
		\begin{compactitem}
			\item Desplazarse a la pestaña de ubicaciones
			\item Pinchar el icono de la derecha de la ubicación a editar
			\item Editar los campos 
		\end{compactitem}
		 \\ \midrule
		\textbf{Resultado esperado} & 
		Se muestra la ubicación con los nuevos datos introducidos
		\\ \midrule
		\textbf{Resultado obtenido} & La ubicación muestra los datos actualizados correctamente\\ \midrule
	\end{tabular}
	\caption{Caso de prueba 25 - Editar ubicaciones}
\end{table}

%Gestion de calidades -----------------------------
\begin{table}[h]
	\centering
	\label{tabla:prueba26}
	\begin{tabular}{@{}
		>{\columncolor[HTML]{FFFFFF}}p {.25\textwidth} p {.75\textwidth}@{}}
		\toprule
		\textbf{Caso de prueba 26}   & \textbf{Visualizar calidades} \\ \midrule
		\textbf{Descripción}     & Comprobar la correcta visualización de las calidades \\ \midrule
		\textbf{Prerrequisitos}	&  Loguearse como gestor de inventario \\ \midrule
		\textbf{Pasos}  & 
		\begin{compactitem}
			\item  Desplazarse a la pestaña de centros
		\end{compactitem}
		 \\ \midrule
		\textbf{Resultado esperado} & Ver la lista de centros
		\\ \midrule
		\textbf{Resultado obtenido} & Visualizada lista de centros\\ \midrule
	\end{tabular}
	\caption{Caso de prueba 26 - Visualizar calidades}
\end{table}

\begin{table}[h]
	\centering
	\label{tabla:prueba27}
	\begin{tabular}{@{}
		>{\columncolor[HTML]{FFFFFF}}p {.25\textwidth} p {.75\textwidth}@{}}
		\toprule
		\textbf{Caso de prueba 27}   & \textbf{Filtrar calidades} \\ \midrule
		\textbf{Descripción}	&  Comprobar el sistema de filtrado de calidades \\ \midrule
		\textbf{Prerrequisitos} & Loguearse como gestor de inventario\\ \midrule
		\textbf{Pasos}  & 
		\begin{compactitem}
			\item Desplazarse a la ventana de calidades desde el menu
			\item Establecer parámetros en la columna de filtros
		\end{compactitem}
		 \\ \midrule
		\textbf{Resultado esperado} & 
		La lista de calidades solo muestra las filas que cumplen los parámetros indicados
		\\ \midrule
		\textbf{Resultado obtenido} & Solo se pueden ver calidades con los parámetros filtrados. \\ \midrule
	\end{tabular}
	\caption{Caso de prueba 27 - Filtrar calidades}
\end{table}

\begin{table}[h]
	\centering
	\label{tabla:prueba28}
	\begin{tabular}{@{}
		>{\columncolor[HTML]{FFFFFF}}p {.25\textwidth} p {.75\textwidth}@{}}
		\toprule
		\textbf{Caso de prueba 28}   & \textbf{Insertar calidades} \\ \midrule
		\textbf{Descripción}	&  Comprobar la correcta inserción de la calidad \\ \midrule
		\textbf{Prerrequisitos}   & Loguearse como gestor de inventario \\ \midrule
		\textbf{Pasos}  & 
		\begin{compactitem}
			\item Desplazarse a la pestaña de centros
			\item Pinchar el botón de Nuevo centro
			\item Insertar los campos
			\item Pulsar en confirmar  
		\end{compactitem}
		 \\ \midrule
		\textbf{Resultado esperado} & 
		Se muestra el nuevo centro en la lista de centros
		\\ \midrule
		\textbf{Resultado obtenido} & la calidad se muestra correctamente \\ \midrule
	\end{tabular}
	\caption{Caso de prueba 28 - Insertar calidades}
\end{table}

\begin{table}[h]
	\centering
	\label{tabla:prueba29}
	\begin{tabular}{@{}
		>{\columncolor[HTML]{FFFFFF}}p {.25\textwidth} p {.75\textwidth}@{}}
		\toprule
		\textbf{Caso de prueba 29}   & \textbf{Editar calidades} \\ \midrule
		\textbf{Descripción}	& Comprobar la correcta actualización de la calidad \\ \midrule
		\textbf{Prerrequisitos} & Loguearse como gestor de inventario \\ \midrule
		\textbf{Pasos}  & 
		\begin{compactitem}
			\item Desplazarse a la pestaña de centros
			\item Pinchar el icono de la derecha de la calidad a editar
			\item Editar los campos 
		\end{compactitem}
		 \\ \midrule
		\textbf{Resultado esperado} & 
		Se muestra la calidad con los nuevos datos introducidos
		\\ \midrule
		\textbf{Resultado obtenido} & la calidad muestra los datos actualizados correctamente\\ \midrule
	\end{tabular}
	\caption{Caso de prueba 29 - Editar calidades}
\end{table}

%Gestion de marcas -----------------------------
\begin{table}[h]
	\centering
	\label{tabla:prueba30}
	\begin{tabular}{@{}
		>{\columncolor[HTML]{FFFFFF}}p {.25\textwidth} p {.75\textwidth}@{}}
		\toprule
		\textbf{Caso de prueba 30}   & \textbf{Visualizar marcas} \\ \midrule
		\textbf{Descripción}     & Comprobar la correcta visualización de las marcass \\ \midrule
		\textbf{Prerrequisitos}	&  Loguearse como gestor de inventario \\ \midrule
		\textbf{Pasos}  & 
		\begin{compactitem}
			\item  Desplazarse a la pestaña de marcas
		\end{compactitem}
		 \\ \midrule
		\textbf{Resultado esperado} & Ver la lista de marcas
		\\ \midrule
		\textbf{Resultado obtenido} & Visualizada lista de marcas\\ \midrule
	\end{tabular}
	\caption{Caso de prueba 30 - Visualizar marcas}
\end{table}

\begin{table}[h]
	\centering
	\label{tabla:prueba31}
	\begin{tabular}{@{}
		>{\columncolor[HTML]{FFFFFF}}p {.25\textwidth} p {.75\textwidth}@{}}
		\toprule
		\textbf{Caso de prueba 31}   & \textbf{Filtrar marcas} \\ \midrule
		\textbf{Descripción}	&  Comprobar el sistema de filtrado de marcas \\ \midrule
		\textbf{Prerrequisitos} & Loguearse como gestor de inventario\\ \midrule
		\textbf{Pasos}  & 
		\begin{compactitem}
			\item Desplazarse a la ventana de marcas desde el menu
			\item Establecer parámetros en la columna de filtros
		\end{compactitem}
		 \\ \midrule
		\textbf{Resultado esperado} & 
		La lista de marcas solo muestra las filas que cumplen los parámetros indicados
		\\ \midrule
		\textbf{Resultado obtenido} & Solo se pueden ver marcas con los parámetros filtrados. \\ \midrule
	\end{tabular}
	\caption{Caso de prueba 31 - Filtrar marcas}
\end{table}

\begin{table}[h]
	\centering
	\label{tabla:prueba32}
	\begin{tabular}{@{}
		>{\columncolor[HTML]{FFFFFF}}p {.25\textwidth} p {.75\textwidth}@{}}
		\toprule
		\textbf{Caso de prueba 32}   & \textbf{Insertar marcas} \\ \midrule
		\textbf{Descripción}	&  Comprobar la correcta inserción de la marca \\ \midrule
		\textbf{Prerrequisitos}   & Loguearse como gestor de inventario \\ \midrule
		\textbf{Pasos}  & 
		\begin{compactitem}
			\item Desplazarse a la pestaña de marcas
			\item Pinchar el botón de Nuevo marca
			\item Insertar los campos
			\item Pulsar en confirmar  
		\end{compactitem}
		 \\ \midrule
		\textbf{Resultado esperado} & 
		Se muestra la nuevo marca en la lista de marcas
		\\ \midrule
		\textbf{Resultado obtenido} & la marca se muestra correctamente \\ \midrule
	\end{tabular}
	\caption{Caso de prueba 32 - Insertar marcas}
\end{table}

\begin{table}[h]
	\centering
	\label{tabla:prueba33}
	\begin{tabular}{@{}
		>{\columncolor[HTML]{FFFFFF}}p {.25\textwidth} p {.75\textwidth}@{}}
		\toprule
		\textbf{Caso de prueba 33}   & \textbf{Editar marcas} \\ \midrule
		\textbf{Descripción}	& Comprobar la correcta actualización de la marca \\ \midrule
		\textbf{Prerrequisitos} & Loguearse como gestor de inventario \\ \midrule
		\textbf{Pasos}  & 
		\begin{compactitem}
			\item Desplazarse a la pestaña de marcas
			\item Pinchar el icono de la derecha de la marca a editar
			\item Editar los campos 
		\end{compactitem}
		 \\ \midrule
		\textbf{Resultado esperado} & 
		Se muestra la marca con los nuevos datos introducidos
		\\ \midrule
		\textbf{Resultado obtenido} & la marca muestra los datos actualizados correctamente\\ \midrule
		\\ \midrule
		\textbf{Resultado obtenido} & La aplicación volverá a mostrar la lista con la marca editada. \\ \midrule
	\end{tabular}
	\caption{Caso de prueba 33 - Editar marcas}
\end{table}

%Gestion de proveedores -----------------------------
\begin{table}[h]
	\centering
	\label{tabla:prueba34}
	\begin{tabular}{@{}
		>{\columncolor[HTML]{FFFFFF}}p {.25\textwidth} p {.75\textwidth}@{}}
		\toprule
		\textbf{Caso de prueba 34}   & \textbf{Visualizar proveedores} \\ \midrule
		\textbf{Descripción}     & Comprobar la correcta visualización de los proveedores \\ \midrule
		\textbf{Prerrequisitos}	&  Loguearse como gestor de inventario \\ \midrule
		\textbf{Pasos}  & 
		\begin{compactitem}
			\item  Desplazarse a la pestaña de proveedores
		\end{compactitem}
		 \\ \midrule
		\textbf{Resultado esperado} & Ver la lista de proveedores
		\\ \midrule
		\textbf{Resultado obtenido} & Visualizada lista de proveedores\\ \midrule
	\end{tabular}
	\caption{Caso de prueba 34 - Visualizar proveedores}
\end{table}

\begin{table}[h]
	\centering
	\label{tabla:prueba35}
	\begin{tabular}{@{}
		>{\columncolor[HTML]{FFFFFF}}p {.25\textwidth} p {.75\textwidth}@{}}
		\toprule
		\textbf{Caso de prueba 35}   & \textbf{Filtrar proveedores} \\ \midrule
		\textbf{Descripción}	&  Comprobar el sistema de filtrado de proveedores \\ \midrule
		\textbf{Prerrequisitos} & Loguearse como gestor de inventario\\ \midrule
		\textbf{Pasos}  & 
		\begin{compactitem}
			\item Desplazarse a la ventana de proveedores desde el menu
			\item Establecer parámetros en la columna de filtros
		\end{compactitem}
		 \\ \midrule
		\textbf{Resultado esperado} & 
		La lista de proveedores solo muestra las filas que cumplen los parámetros indicados
		\\ \midrule
		\textbf{Resultado obtenido} & Solo se pueden ver proveedores con los parámetros filtrados. \\ \midrule
	\end{tabular}
	\caption{Caso de prueba 35 - Filtrar proveedores}
\end{table}

\begin{table}[h]
	\centering
	\label{tabla:prueba36}
	\begin{tabular}{@{}
		>{\columncolor[HTML]{FFFFFF}}p {.25\textwidth} p {.75\textwidth}@{}}
		\toprule
		\textbf{Caso de prueba 36}   & \textbf{Insertar proveedores} \\ \midrule
		\textbf{Descripción}	&  Comprobar la correcta inserción del proveedor \\ \midrule
		\textbf{Prerrequisitos}   & Loguearse como gestor de inventario \\ \midrule
		\textbf{Pasos}  & 
		\begin{compactitem}
			\item Desplazarse a la pestaña de proveedores
			\item Pinchar el botón de Nuevo proveedor
			\item Insertar los campos
			\item Pulsar en confirmar  
		\end{compactitem}
		 \\ \midrule
		\textbf{Resultado esperado} & 
		Se muestra el nuevo proveedor en la lista de proveedores
		\\ \midrule
		\textbf{Resultado obtenido} & El proveedor se muestra correctamente \\ \midrule
	\end{tabular}
	\caption{Caso de prueba 36 - Insertar proveedores}
\end{table}

\begin{table}[h]
	\centering
	\label{tabla:prueba37}
	\begin{tabular}{@{}
		>{\columncolor[HTML]{FFFFFF}}p {.25\textwidth} p {.75\textwidth}@{}}
		\toprule
		\textbf{Caso de prueba 37}   & \textbf{Editar proveedores} \\ \midrule
		\textbf{Descripción}	& Comprobar la correcta actualización del proveedor \\ \midrule
		\textbf{Prerrequisitos} & Loguearse como gestor de inventario \\ \midrule
		\textbf{Pasos}  & 
		\begin{compactitem}
			\item Desplazarse a la pestaña de proveedores
			\item Pinchar el icono de la derecha del proveedor a editar
			\item Editar los campos 
		\end{compactitem}
		 \\ \midrule
		\textbf{Resultado esperado} & 
		Se muestra el proveedor con los nuevos datos introducidos
		\\ \midrule
		\textbf{Resultado obtenido} & El proveedor muestra los datos actualizados correctamente\\ \midrule
	\end{tabular}
	\caption{Caso de prueba 37 - Editar proveedores}
\end{table}

Por otro lado, también se han desarrollado unas pruebas \textbf{alpha}. Estas pruebas se realizaron de manera telemática junto con los miembros de laboratorio que serán los usuarios de la aplicación.
\apendice{Documentación de usuario}

\section{Introducción}

En este apartado se realizará una guía detallada para que el usuario sepa cómo utilizar la aplicación.

\section{Requisitos de usuarios}

Al ser una aplicación web, el cliente únicamente necesita tener instalado un navegador web que soporte \textit{javascript}, \textit{jQuery} y \textit{hojas de estilo CSS}. 

La aplicación funciona en cualquier navegador, pero se recomienda su uso en Google Chrome.

\section{Instalación}

El cliente no necesita realizar ninguna instalación, únicamente ha de acceder a la dirección web en la que se aloja la aplicación.

\section{Manual del usuario}

En este apartado se explicará cada una de las acciones que puede realizar el cliente sobre nuestra aplicación. Lo dividiremos en diferentes acciones según cada uno de los contextos de la aplicación.

\begin{itemize}
	\item Login
	\item Página de administración (Administrador)
	\item Página de gestión (Gestor de inventario)
\end{itemize}

\subsection{Login}

La primera página que nos encontramos es la página de login \ref{fig:pantallaLogin}. Desde aquí el usuario puede acceder a su página de usuario correspondiente en función de su rol (administrador, gestor de inventario o usuario).

\imagen{pantallaLogin}{Pantalla de login}{1} 

\subsection{Página de administración}

Si el usuario tiene rol de administrador accede a la página de administración. En ella se puede mover por las distintas pestañas que pertenecen a los centros \ref{fig:administracion-centros}, departamentos \ref{fig:administracion-dptos}, áreas \ref{fig:administracion-areas} o usuarios.

\imagen{administracion-centros}{Pantalla administración - Pestaña centros}{1} 
\imagen{administracion-dptos}{Pantalla administración - Pestaña departamentos}{1} 
\imagen{administracion-areas}{Pantalla administración - Pestaña áreas}{1} 
\imagen{administracion-usuarios}{Pantalla administración - Pestaña usuarios}{1} 

Desde cualquiera de estas pestañas el usuario puede hacer varias acciones, puede insertar un elemento o editar un elemento. A continuación, mostraremos unos ejemplos, añadiendo un área y editando un usuario. El resto de procesos no se mostrarán ya que son exactamente iguales, con la diferencia de algún parámetro.

\subsubsection{Insertar}

Para insertar un elemento solo hay que pulsar en el botón verde de abajo a la derecha donde dice Nuevo área, por ejemplo. Hecho esto se creará una nueva línea en la lista de áreas en la que podremos añadir los datos \ref{fig:administracion-insertar}. Con los datos ya introducidos solo hay que pulsar en confirmar (verde) o en cancelar (rosa), en los botones de abajo.

\imagen{administracion-insertar}{Pantalla administración - Insertar área}{0.8}

\subsubsection{Editar}

Para editar un elemento hay que pulsar en el botón rosa con un lapicero en la columna de la derecha de la fila que deseamos editar. Hecho esto se habilitará la edición de los valores de la fila y podremos editar los datos \ref{fig:administracion-editar}. Con los datos ya introducidos solo hay que pulsar en confirmar (verde) o en cancelar (rosa), en los botones la derecha donde antes pulsamos en editar.

En el caso de usuario aparece además un botón azul, que sirve para reiniciar la contraseña del usuario.

\imagen{administracion-editar}{Pantalla administración - Editar usuario}{0.8}

\subsubsection{Cerrar sesión}

Si el usuario desea cerrar sesión solamente ha de pinchar en su nombre arriba a la derecha para mostrar el menú, donde le aparecerá la opción de cerrar sesión \ref{fig:administracion-cerrarSesion} 

\imagen{administracion-cerrarSesion}{Pantalla administración - Cerrar sesión}{0.4}

\subsection{Página de gestión}

Si el usuario tiene un rol de gestión de inventario accede directamente a la página de gestión \ref{fig:gestion-existencias}. Desde aquí el usuario puede moverse al resto de páginas desde el menú, pulsando arriba a la derecha sobre su nombre \ref{fig:gestion-menu}

\imagen{gestion-existencias}{Página de existencias}{1}
\imagen{gestion-menu}{Menú}{0.3}

Si queremos ver más información sobre alguna ficha en cuestión, solo tenemos que pinchar sobre su fila y se abrirá una pequeña ventana con información complementaria \ref{fig:gestion-existencias-info}.

\imagen{gestion-existencias-info}{Extra info existencias}{0.8}

Si queremos añadir una entrada o unas salida de una ficha de producto, solo hay que pulsar en los iconos de la derecha de su fila, en el icono de un + (botón azul) para añadir una entrada, y en el icono de un - (botón rosa) para añadir una salida.

Una vez pulsado uno de los dos botones se nos abrirá una pequeña ventana modal que nos pide el número de unidades, la fecha y si queremos añadir alguna nota \ref{fig:gestion-entrada} \ref{fig:gestion-salida}.  

\imagen{gestion-entrada}{Añadir una entrada}{0.5}

\imagen{gestion-salida}{Añadir una salida}{0.5}

\subsubsection{Insertar / Editar}

Si queremos insertar un elemento solamente hay que pulsar en el botón verde de abajo a la derecha y se abrirá un modal \ref{fig:gestion-insertar-editar} con todos los datos a rellenar. Si lo queremos es editar, hay que pulsar en el botón rosa con icono de un lapicero en la columna de la derecha de la fila del elemento que queremos editar. Al pinchar se abrirá la misma ventana modal que al insertar \ref{fig:gestion-insertar-editar}, pero con los datos ya rellenados. Ahí podemos editar los datos necesarios y pulsar en confirmar.  
Para insertar o editar un elemento en cualquiera de las páginas el proceso es muy parecido, únicamente cambiando los datos que aparecen. 

\imagen{gestion-insertar-editar}{Modal insertar nueva ficha}{0.8}

\subsubsection{Página productos}

Si se accede a la página de productos se muestra la lista con los productos que existen en el sistema \ref{fig:gestion-productos}.

\imagen{gestion-productos}{Página de productos}{1}

Si queremos ver más información sobre algún producto en cuestión, solo tenemos que pinchar sobre su fila y se abrirá una pequeña ventana con información complementaria \ref{fig:gestion-productos-info}.

\imagen{gestion-productos-info}{Extra info productos}{0.8}

\subsubsection{Página ubicaciones}

Si se accede a la página de productos se muestra la lista con las ubicaciones que existen en el sistema \ref{fig:gestion-ubicaciones}.

\imagen{gestion-ubicaciones}{Página de ubicaciones}{1}

\subsubsection{Página calidades}

Si se accede a la página de productos se muestra la lista con las ubicaciones que existen en el sistema \ref{fig:gestion-ubicaciones}.

\imagen{gestion-calidades}{Página de calidades}{1}

\subsubsection{Página marcas}

Si se accede a la página de productos se muestra la lista con las marcas que existen en el sistema \ref{fig:gestion-marcas}.

\imagen{gestion-marcas}{Página de marcas}{1}

Si queremos ver más información sobre algúna marca en cuestión, solo tenemos que pinchar sobre su fila y se abrirá una pequeña ventana con información complementaria \ref{fig:gestion-marcas-info}.

\imagen{gestion-marcas-info}{Extra info marcas}{0.8}

\subsubsection{Página proveedores}

Si se accede a la página de productos se muestra la lista con los proveedores que existen en el sistema \ref{fig:gestion-proveedores}.

\imagen{gestion-proveedores}{Página de proveedores}{0.9}

Si queremos ver más información sobre algún proveedor en cuestión, solo tenemos que pinchar sobre su fila y se abrirá una pequeña ventana con información complementaria \ref{fig:gestion-proveedores-info}.

\imagen{gestion-proveedores-info}{Extra info proveedores}{0.8}


\bibliographystyle{plain}
\bibliography{bibliografiaAnexos}

\end{document}
